% ****** Start of file main.tex ******
%
\documentclass[%
 amsmath,amssymb,
 aps, onecolumn
]{revtex4-2}

\usepackage{graphicx}% Include figure files
\usepackage{dcolumn}% Align table columns on decimal point
\usepackage{bm}% bold math

\usepackage{hyperref}% required for \href and \nolinkurl in bibliography
\usepackage{aas_macros}% AAS journal abbreviations used in references.bib

\newcommand{\ResultDir}{../output/DR1_v4}   % output directory for figures & numbers


%% ---------------------------------------------------------------
%% Synced values — updated automatically by /sync-numbers
%% Do not edit values by hand; run /sync-numbers after a DR1 re-fit.
%% ---------------------------------------------------------------
% config.json
\newcommand{\Ntrain}{5{,}050}           % n_training
\newcommand{\Ntotal}{10{,}200}          % n_total_fits
% selection_ellipse.json
\newcommand{\GmmW}{0.935}               % weight
\newcommand{\GmmMeanX}{0.171}           % mean[0]
\newcommand{\GmmMeanY}{-20.409}         % mean[1]
\newcommand{\GmmSemiA}{0.064}           % semi_axes[0]
\newcommand{\GmmSemiB}{0.729}           % semi_axes[1]
\newcommand{\GmmAngle}{98.7}            % angle_deg
% select_v2_mle.json
\newcommand{\MleSlope}{-6.8346}         % slope
\newcommand{\MleIntercept}{-19.2215}    % intercept.1
\newcommand{\MleSigIntX}{6.132 \times 10^{-2}}  % sigma_int_x
\newcommand{\MleSigIntY}{0.0991}        % sigma_int_y
% select_v2_fiducial.json
\newcommand{\FidHatyMin}{-21.85}        % haty_min
\newcommand{\FidHatyMax}{-19.00}        % haty_max
\newcommand{\FidSlopePlane}{-6.535}     % slope_plane
\newcommand{\FidCOne}{-20.565}          % intercept_plane
\newcommand{\FidCTwo}{-18.021}          % intercept_plane2
\newcommand{\FidNSigPerp}{3.0}          % n_sigma_perp
\newcommand{\FidZMin}{0.03}             % z_obs_min
\newcommand{\FidZMax}{0.08}             % z_obs_max
% tophat_?.csv — slope (α)
\newcommand{\SlopeMean}{-6.691}         % mean
\newcommand{\SlopeStd}{0.073}           % std
\newcommand{\SlopeQlo}{-6.816}          % 5%
\newcommand{\SlopeQmed}{-6.689}         % 50%
\newcommand{\SlopeQhi}{-6.573}          % 95%
\newcommand{\SlopeEss}{1178}            % ESS_bulk (stansummary)
\newcommand{\SlopeRhat}{1.0}            % R_hat (stansummary)
% tophat_?.csv — intercept (β)
\newcommand{\IntMean}{-19.245}          % mean
\newcommand{\IntStd}{0.015}             % std
\newcommand{\IntQlo}{-19.270}           % 5%
\newcommand{\IntQmed}{-19.246}          % 50%
\newcommand{\IntQhi}{-19.219}           % 95%
\newcommand{\IntEss}{1301}              % ESS_bulk (stansummary)
\newcommand{\IntRhat}{1.0}              % R_hat (stansummary)
% tophat_?.csv — sigma_int_x
\newcommand{\SigXMean}{0.013}           % mean
\newcommand{\SigXStd}{0.006}            % std
\newcommand{\SigXQlo}{0.003}            % 5%
\newcommand{\SigXQmed}{0.014}           % 50%
\newcommand{\SigXQhi}{0.022}            % 95%
\newcommand{\SigXEss}{875}              % ESS_bulk (stansummary)
\newcommand{\SigXRhat}{1.0}             % R_hat (stansummary)
% tophat_?.csv — sigma_int_y
\newcommand{\SigYMean}{0.438}           % mean
\newcommand{\SigYStd}{0.011}            % std
\newcommand{\SigYQlo}{0.420}            % 5%
\newcommand{\SigYQmed}{0.439}           % 50%
\newcommand{\SigYQhi}{0.456}            % 95%
\newcommand{\SigYEss}{1414}             % ESS_bulk (stansummary)
\newcommand{\SigYRhat}{1.0}             % R_hat (stansummary)

\begin{document}

\preprint{APS/123-QED}

\title{Model-Based Calibration of the Tully-Fisher Relation for DESI DR2 Peculiar Velocities}

%\author[0000-0001-6315-8743]{Alex G.\ Kim}
\author{Alex G.\ Kim}
\email{agkim@lbl.gov}
%\homepage[]{Your web page}
%\thanks{}
%\altaffiliation{}
\affiliation{Physics Division, Lawrence Berkeley National Laboratory\\
1 Cyclotron Road, Berkeley, CA 94720 USA}


\collaboration{DESI Collaboration}%\noaffiliation

\date{\today}

\begin{abstract}
[DRAFT]
We present a model-based calibration of the Tully-Fisher relation (TFR) using
the full spiral-galaxy sample from the DESI Data Release 2 (DR2) Peculiar
Velocity survey.  The DR2 catalog contains $\Ntotal$ spiral galaxies drawn from
the Siena Galaxy Atlas.  We define the TFR phase-space variables
$x = \log_{10}(V_\text{rot}/100~\text{km\,s}^{-1})$ and $y = M_R$
(the $r$-band isophotal absolute magnitude at $SB=26$~mag\,arcsec$^{-2}$),
and fit a two-component Gaussian mixture model (GMM) to the observed
$(x,\,y)$ distribution to characterize the spiral population.
The dominant GMM component (weight $w=\GmmW$) defines the orientation of
oblique selection cuts in phase space; galaxies passing these cuts together
with a redshift window $0.03 < z < 0.1$ constitute the training sample of
$\Ntrain$ galaxies.  We calibrate the TFR by fitting a Stan tophat model to
this sample via Hamiltonian Monte Carlo, obtaining slope
$\alpha = \SlopeMean \pm \SlopeStd$, intercept $\beta = \IntMean \pm \IntStd$,
intrinsic scatter along the velocity axis $\sigma_{\mathrm{int},x} = \SigXMean \pm \SigXStd$,
and intrinsic scatter along the magnitude axis $\sigma_{\mathrm{int},y} = \SigYMean \pm \SigYStd$.
The calibrated model is applied to the full DR2 spiral sample to produce a
peculiar velocity catalog (\texttt{tophat\_catalog.fits}) and a full posterior
predictive covariance matrix (\texttt{tophat\_cov.fits}).
\end{abstract}

\maketitle


\section{Introduction}
\label{sec:intro}
Peculiar velocities --- the deviations of galaxy redshifts from pure Hubble flow ---
encode the gravitational influence of large-scale structure and provide an independent
probe of the growth rate of density perturbations, bulk flows, and the Hubble constant
\citep[e.g.][]{1997AJ....113...53G, 2008ApJ...676..184T, 2019MNRAS.487.2061H}.
The Tully-Fisher relation \citep[TFR;][]{Tully1977} links the intrinsic luminosity of
a spiral galaxy to its rotation velocity.
Because more luminous spirals rotate faster, a consequence of the dark-matter halos
that govern galaxy dynamics \citep{1998MNRAS.295..319M},
the TFR provides a distance estimate independent of redshift: the luminosity is inferred
from the rotation velocity, and the distance from the luminosity--flux ratio.
Peculiar velocity surveys based on the TFR have been carried out by a number of groups
\citep[e.g.][]{2000ApJ...529..698S, 2008ApJ...676..184T, 2019MNRAS.487.2061H,
2020MNRAS.497.1275S, 2020ApJ...902..145K}.

The Dark Energy Spectroscopic Instrument \citep[DESI;][]{cite_desi} is now delivering
the largest spectroscopic peculiar-velocity dataset to date.
The DESI Peculiar Velocity survey \citep{Saulder2023} targets spiral galaxies,
early-type galaxies (Fundamental Plane), and supernova hosts.
A first TFR catalog of 1,106 spirals from the DESI Early Data Release
\citep[EDR;][]{DESI_EDR} was presented in \citet{2024Douglass}, with particular depth
in the Coma cluster region.
Data Release 1 (DR2) extends the spiral sample to $\Ntotal$ galaxies spanning the full DESI
footprint, forming the dataset analyzed here.

Traditional TFR calibrations fit a line to the data after applying iterative
bias corrections for selection effects and measurement errors
\citep{2006ApJ...653..861M, 2012ApJ...749...78T, 1997AJ....113...53G}.
We instead adopt a forward-modeling approach: we write down a probabilistic galaxy
model that contains the TFR as its core relationship, and calibrate it by maximizing
the likelihood of the observations.
This has several advantages over the traditional approach.
Measurement uncertainties in both the rotation velocity and magnitude, including the
intrinsic scatter of the TFR, enter the likelihood directly without the need for
separate bias corrections.
The selection function is folded into the likelihood, so no iterative reweighting of
the sample is required.
The TFR slope, intercept, and intrinsic scatter parameters are all determined
in a single joint fit.
This modeling approach does require specifying the intrinsic distribution of galaxy
rotation velocities; we address this by fitting a Gaussian mixture model to the
observed phase space and adopting a uniform (tophat) prior within the selected region
(\S\ref{sec:model}--\ref{sec:selection}).

In this paper we present a model-based calibration of the TFR using the DESI DR2
spiral sample.
We describe the sample selection criteria (\S\ref{sec:sampleselection}),
the Stan tophat model and its likelihood (\S\ref{sec:model}),
the GMM-based phase-space selection function (\S\ref{sec:selection}),
the MCMC fit results and diagnostics (\S\ref{sec:results}),
and the output peculiar velocity catalog (\S\ref{sec:catalog}).


\section{Data}
\label{sec:data}
The DESI DR2 Peculiar Velocity spiral sample contains $\Ntotal$ galaxies drawn from the Siena Galaxy Atlas \citep[SGA;][]{SGA}, which is based on the DESI Legacy Imaging Surveys DR9 data \citep{DESI_Imaging} and covers areas north of $-30^\circ$ declination. The SGA includes galaxies with diameters larger than 20'' at the 25th
magnitude isophote, denoted as $D(25)$. To be part of the TFR sample, a galaxy must have a Sersic index less than 2, indicating it is likely a spiral galaxy, and an inclination angle greater than $25^\circ$, ensuring a significant portion of the galaxy's rotational velocity is visible and the galaxy's major axis is properly aligned.

SGA provides galaxy shape and magnitude information
at the $r$-band $SB=26$~$\text{mag}\,\text{arcsec}^{-2}$ isophote, in the form of the photometric axis ratio $b/a$,
the position angle PA of the semi-major axis,
and the radius $R_{26}$. 

The DESI measurements are taken at
the galactic center and on both sides along the major axis oriented according to the SGA PA.
Rotational velocities were measured at $0.4R_{26}$,
the distance at which reliable redshifts could consistently be measured.  The redshifts measured at different positions
are combined to obtain a redshift $z$ and measured
rotational velocity $V_\text{obs}$ for each galaxy
\citep{2024Douglass}.

SGA and DESI data are combined to infer the rotation
velocity $V_\text{rot}$.
The galaxy inclination angle $i$ is derived from
\begin{equation}\label{eqn:BAtoi}
    \cos^2 i = \frac{(b/a)^2 - q_0^2}{1 - q_0^2},
\end{equation}
where $q_0 = 0.2$ \citep{Tully2000}.
This is in
turn used to convert the observed line-of-sight velocity to the
rotation velocity as
\begin{equation}\label{eqn:Vcorrected}
    V_{\text{rot}} = \frac{V_{\text{obs}}}{\sin i}.
\end{equation}

The analysis variables used throughout this paper are
\begin{align}
    x &= \log_{10}\!\left(\frac{V_\text{rot}}{100~\text{km\,s}^{-1}}\right), \\
    y &= M_R,
\end{align}
where $M_R$ is the $r$-band SB26 absolute magnitude computed from
the apparent magnitude and the luminosity distance at the observed redshift.

The intrinsic galaxy distribution is redshift-dependent, as illustrated in
Figure~\ref{fig:redshift_grid}, which shows the average redshift of DR2
spiral galaxies in bins of rotation velocity and absolute magnitude.
Intrinsically faint galaxies are preferentially observed
at low redshift, reflecting a genuine property of the local universe rather
than an observational bias: the SGA catalog is selected on angular size
($D(25) > 20''$) and is therefore not subject to apparent-magnitude
selection effects.
However, the low-redshift sample is more heavily contaminated by dwarf
galaxies, which do not follow the TFR and are thus poorly suited for
calibration.
We therefore apply a minimum redshift cut for the TFR calibration sample,
and caution that TFR-inferred distances for low-redshift galaxies may be
more susceptible to contamination by dwarfs that do not follow the TFR.
This redshift dependence of the galaxy population also propagates into
redshift-dependent covariance in the TFR-inferred peculiar velocities.

\begin{figure}
    \centering
    \includegraphics[width=1\linewidth]{\ResultDir/redshift_grid_tophat.png}
    \caption{Average redshift of DR2 spiral galaxies in bins of
    $\log_{10}(V_\text{rot}/100~\text{km\,s}^{-1})$ and $M_R$.
    Intrinsically faint, slowly rotating galaxies are preferentially
    found at low redshift, reflecting the local galaxy population rather
    than an observational selection effect.}
    \label{fig:redshift_grid}
\end{figure}

\section*{H$\alpha$ Redshift Range in DECam $r$-band}

The H$\alpha$ emission line has a rest-frame wavelength of $\lambda_{\rm H\alpha} = 656.3$ nm. 
The DECam $r$-band spans approximately $\lambda_{\rm blue} = 562$ nm to $\lambda_{\rm red} = 698$ nm 
at the 50\% transmission points.

The observed wavelength of a spectral line is related to its rest-frame wavelength by:
\begin{equation}
    \lambda_{\rm obs} = \lambda_{\rm emit} \,(1 + z),
\end{equation}
so the corresponding redshift is:
\begin{equation}
    z = \frac{\lambda_{\rm obs}}{\lambda_{\rm emit}} - 1.
\end{equation}

At the blue edge of the $r$-band:
\begin{equation}
    z_{\rm blue} = \frac{562}{656.3} - 1 \approx -0.14,
\end{equation}
which corresponds to a blueshift and is not physical in a cosmological context.

At the red edge of the $r$-band:
\begin{equation}
    z_{\rm red} = \frac{698}{656.3} - 1 \approx 0.064.
\end{equation}

Since H$\alpha$ at rest ($z = 0$) falls at 656.3 nm, which is already within the $r$-band, 
the H$\alpha$ line is observable in the DECam $r$-band over the redshift range:
\begin{equation}
    0 \lesssim z \lesssim 0.06.
\end{equation}

\section{Sample Selection}
\label{sec:sampleselection}
Not all galaxies in the DR2 spiral catalog are suitable for TFR calibration.
Dwarf and irregular galaxies at the faint end of the luminosity function do
not follow the same TFR as normal spirals, while extreme outliers at either
end of the phase-space distribution can disproportionately influence the fit.
We therefore restrict the analysis to a region $\mathcal{S}$ in the observed
$(\hat{x},\hat{y})$ plane that selects the well-behaved spiral population.
The selection process has two phases, with the second region a refinement of the first.

The region $\mathcal{S}$ is defined by two types of cut.
First, upper and lower limits on the observed absolute magnitude $\hat{y}$
exclude the faintest galaxies, which depart from the TFR, and the
brightest galaxies, which are rare and susceptible to outlier contamination.
Second, two diagonal bands oriented approximately parallel to the TFR --- and therefore
perpendicular to the direction of TFR scatter --- bound the sample on either
side of the main locus.
Cutting parallel to the relation rather than in $\hat{y}$ or $\hat{x}$ alone
efficiently encloses the TFR population across the full range of rotation
velocities and avoids sculpting the scatter distribution asymmetrically.
Together the four boundaries define a parallelogram-shaped region,
\begin{equation}
(\hat{x},\hat{y})\in\mathcal{S}
\;\Longleftrightarrow\;
\hat{y}_{\min} \le \hat{y} \le \hat{y}_{\max}
\;\text{ and }\;
\bar{s}\,\hat{x}+\bar{c}_1 \le \hat{y} \le \bar{s}\,\hat{x}+\bar{c}_2,
\label{eqn:selection_region}
\end{equation}
where $\hat{y}_{\min}$ and $\hat{y}_{\max}$ are the magnitude limits,
$\bar{s}$ is the slope of the oblique boundaries,
$\bar{c}_1 < \bar{c}_2$ are their intercepts, and $d\equiv \bar{c}_2-\bar{c}_1$ is the
perpendicular band width.


The TFR is unknown a priori, so an initial estimate is used to define
the slope of the distribution.  Variations in the selected slope perturb
the galaxies that enter the analysis, but otherwise do not bias the fits:
our pipeline recovers the input slopes of mock surveys generated with a consistent model, independent of the initial estimate.
The procedures for selecting the selection regions
are presented in \S\ref{sec:training_selection}, 
after the description of the model.

\section{Model}
\label{sec:model}
We construct a probabilistic forward model for the TFR that accounts for
measurement uncertainties in both observables and for sample selection within
a single likelihood.
The full derivation is given in Appendix~\ref{sec:likelihood};
here we state the elements of the model and the expression for the likelihood.

The sample satisfies the selection criteria
$(\hat{x}_i,\hat{y}_i) \in \mathcal{S}$.
The training sample consists of $N$ galaxies.
For a galaxy indexed by $i$, the observables are the measured log-velocity
$\hat{x}_i$ with uncertainty $\sigma_{x,i}$ and the observed absolute
magnitude $\hat{y}_i$ with uncertainty $\sigma_{y,i}$.  Magnitude
and log-velocity uncertainties are treated as Gaussian.

The global TFR parameters are the slope $\alpha$, intercept $\beta$, and intrinsic
scatter amplitudes $\sigma_{\mathrm{int},x}$ and $\sigma_{\mathrm{int},y}$
along the $x$- and $y$-axes respectively.
Each galaxy has a latent on-TFR magnitude $y_{\mathrm{TF},i}$ that
places it on the relation.   The true absolute magnitude and
velocities $x_i$ and $y_i$
have scatter around the TFR,
\begin{align}
y_i \mid y_{\mathrm{TF},i}
    &\sim \mathcal{N}\!\left(y_{\mathrm{TF},i},\,\sigma_{\mathrm{int},y}^2\right), \\
x_i \mid y_{\mathrm{TF},i}
    &\sim \mathcal{N}\!\left(\frac{y_{\mathrm{TF},i}-\beta}{\alpha},\,\sigma_{\mathrm{int},x}^2\right).
\end{align}
Measurements have Gaussian scatter around the true values:
\begin{align}
\hat{x}_i \mid x_i &\sim \mathcal{N}\!\left(x_i,\,\sigma_{x,i}^2\right), \\
\hat{y}_i \mid y_i &\sim \mathcal{N}\!\left(y_i,\,\sigma_{y,i}^2\right).
\end{align}

For a hierarchical model, the underlying distribution of the latent
parameter must be specified. We adopt a uniform distribution,
$y_{\mathrm{TF},i} \sim \mathrm{Uniform}(y_{\min},\,y_{\max})$, with
bounds extending slightly beyond the selection window to accommodate
scatter at the selection boundary,
\begin{align}
y_{\min} &= \hat{y}_{\min} - \Delta y_{\min}, &
y_{\max} &= \hat{y}_{\max} + \Delta y_{\max},
\label{eqn:ymin_ymax}
\end{align}
where we set $\Delta y_{\min}=-0.5$, $\Delta y_{\max}=1.0$.
The choice of latent model matters only for objects whose latent
distribution varies appreciably across the width of the Gaussian
measurement error; for objects with small uncertainties the likelihood
is insensitive to the model choice. To minimize sensitivity to
this choice, we place the selectrion criteria bounds in a region where $p(y_{\mathrm{TF}})$
is expected to have a weak gradient, so that the uniform remains a good
local approximation even for objects with broad likelihoods. In this
regime the normalization $P(S_i = 1 \mid \theta)$ varies only weakly
with small changes to $\Delta y_{\min}$ and $\Delta y_{\max}$ (chosen
to be signifcantly
larger than the absolute magnitude measurement uncertainties), and the
inferred TFR parameters are insensitive to the precise edge placements.
This choice is also supported empirically: in mock-catalog tests,
the uniform prior yielded more stable parameter recovery than a Gaussian prior,
consistent with the relatively flat interior of the galaxy absolute-magnitude distribution.

The above model specifies the likelihood.
Given the independence of the galaxy measurements, the total likelihood can be decomposed
into the product of per-galaxy terms. 
\begin{align}
\mathcal{L}(\theta;\{\hat{x}_i\},\{\hat{y}_i\})
&= \prod_{i=1}^{N} p(\hat{x}_i,\hat{y}_i \mid \theta, (\hat{x}_i,\hat{y}_i)\in\mathcal{S}).
\label{eqn:total_likelihood}
\end{align}
Convolving intrinsic and measurement scatter defines the combined uncertainties
\begin{align}
\sigma_{1,i}^2 &= \sigma_{x,i}^2+\sigma_{\mathrm{int},x}^2, &
\sigma_{2,i}^2 &= \sigma_{y,i}^2+\sigma_{\mathrm{int},y}^2.
\label{eqn:sigma12}
\end{align}

Because the sample is restricted to $(\hat{x}_i,\hat{y}_i)\in\mathcal{S}$,
the per-galaxy density is truncated and renormalized,
\begin{align}
p(\hat{x}_i,\hat{y}_i \mid \theta,(\hat{x}_i,\hat{y}_i)\in\mathcal{S})
&= \frac{\mathcal{L}_i\;
         \mathbf{1}\!\left\{(\hat{x}_i,\hat{y}_i)\in\mathcal{S}\right\}}
        {P(S_i=1\mid\theta)}.
\label{eqn:selected_density}
\end{align}
The numerator is the per-galaxy likelihood if there were no  sample selection,
\begin{align}
\mathcal{L}_i
&= \frac{|\alpha|}{y_{\max}-y_{\min}}\;
   \mathcal{N}\!\left(\hat{y}_i;\,\beta+\alpha\hat{x}_i,\,\sigma_{\mathrm{tot},i}^2\right)
   \left[
     \Phi\!\left(\frac{y_{\max}-\mu_{*,i}}{\sigma_{*,i}}\right)
     -\Phi\!\left(\frac{y_{\min}-\mu_{*,i}}{\sigma_{*,i}}\right)
   \right],
\label{eqn:Li_tophat}
\end{align}
where $\Phi$ is the standard normal CDF and
\begin{align}
\sigma_{\mathrm{tot},i}^2
    &= \alpha^2\sigma_{1,i}^2+\sigma_{2,i}^2, \\
\mu_{*,i}
    &= \frac{(\alpha\hat{x}_i+\beta)\,\sigma_{2,i}^2
             +\hat{y}_i\,\alpha^2\sigma_{1,i}^2}{\sigma_{\mathrm{tot},i}^2}, \\
\sigma_{*,i}^2
    &= \frac{\alpha^2\sigma_{1,i}^2\,\sigma_{2,i}^2}{\sigma_{\mathrm{tot},i}^2}.
\label{eqn:posterior_params}
\end{align}


The denominator, i.e.\ the selection probability, is
derived by transforming the $\hat{x}$--$\hat{y}$ to a new coordinate system where
the parallelogram transforms to an axis-aligned rectangle.  The rectangular region is the intersection of four closed half-planes.  The selection probability of
a half plane is the standard bivariate-normal CDF.  The resulting expression is
\begin{align}
P(S_i=1\mid\theta)
&= \frac{1}{y_{\max}-y_{\min}}
   \int_{y_{\min}}^{y_{\max}} I_i(y_{\mathrm{TF}})\,dy_{\mathrm{TF}},
\label{eqn:Psel_tophat}
\end{align}
where $I_i$ is a combination of four bivariate-normal CDF evaluations,
\begin{align}
I_i(y_{\mathrm{TF}})
&\equiv
  \Phi_2\!\Bigl(h_i(0;y_{\mathrm{TF}}),\;k_i(\hat{y}_{\max};y_{\mathrm{TF}});\;{-\rho_i}\Bigr)
 -\Phi_2\!\Bigl(h_i(d;y_{\mathrm{TF}}),\;k_i(\hat{y}_{\max};y_{\mathrm{TF}});\;{-\rho_i}\Bigr)
\nonumber\\
&\quad
 -\Phi_2\!\Bigl(h_i(0;y_{\mathrm{TF}}),\;k_i(\hat{y}_{\min};y_{\mathrm{TF}});\;{-\rho_i}\Bigr)
 +\Phi_2\!\Bigl(h_i(d;y_{\mathrm{TF}}),\;k_i(\hat{y}_{\min};y_{\mathrm{TF}});\;{-\rho_i}\Bigr),
\label{eqn:Ii_def}
\end{align}
with $\Phi_2(h,k;\varrho)$ the standard bivariate-normal CDF,
$d\equiv \bar{c}_2-\bar{c}_1$, and the standardised bounds
\begin{align}
h_i(t;y_{\mathrm{TF}}) &\equiv \frac{\mu_{u,i}(y_{\mathrm{TF}})-t}{\sigma_{u,i}}, &
k_i(y_\star;y_{\mathrm{TF}}) &\equiv \frac{y_\star - y_{\mathrm{TF}}}{\sigma_{2,i}}.
\label{eqn:ab_bounds}
\end{align}
Here $\mu_{u,i}$, $\sigma_{u,i}$, and $\rho_i$ are the mean, standard deviation,
and correlation of the sheared observable $\hat{u}_i\equiv\hat{y}_i - \bar{s}\hat{x}_i - \bar{c}_1$
conditional on $y_{\mathrm{TF}}$:
\begin{align}
\mu_{u,i}(y_{\mathrm{TF}})
    &= y_{\mathrm{TF}} - \bar{s}\,\frac{y_{\mathrm{TF}}-\beta}{\alpha} - \bar{c}_1, \\
\sigma_{u,i}^2
    &= \sigma_{2,i}^2 + \bar{s}^2\,\sigma_{1,i}^2, \\
\rho_i
    &= \frac{\sigma_{2,i}}{\sigma_{u,i}}.
\label{eqn:shear_params}
\end{align}


\section{Tully-Fisher Relation Fit}
\label{sec:TFRfit}
\subsection{Training Set Selection}
\label{sec:training_selection}
In this section we describe the procedure for determining the parameters
that define the sample selection region $\mathcal{S}$
(Eq.~\ref{eqn:selection_region}), as introduced in
\S\ref{sec:sampleselection}. The procedure has three steps: (i) a
Gaussian Mixture Model (GMM) fit identifies the core spiral population
and defines a preliminary selection parallelogram; (ii) a nominal TFR is
fit to the galaxies within that parallelogram; and (iii) a
residual-profile diagnostic applied to the full catalog refines the
magnitude limits to exclude regions where magnitude deviations are significant.

Velocity--magnitude cuts are designed to exclude outliers and sensitivity
to the modeling of the latent parameter distribution. We apply a
minimum and maximum redshift cuts ($0.03 < z < 0.08$) to reduce sensitivity to
peculiar-velocity contamination and potential contamination.

The primary contaminants are dwarf galaxies, which are known not to
follow the TFR, and bright galaxies lying on the sharp declining tail
of the true luminosity function, where the top-hat model may not yield
accurate measurement probabilities. To identify magnitude limits that
excise these populations, we apply the nominal TFR to all galaxies and
examine the residuals between predicted and observed magnitudes as a
function of absolute magnitude. Because the nominal TFR need not be
accurate, a smooth linear trend in the residuals is expected and
acceptable; we instead look for deviations from this trend, which signal
contamination by populations that do not follow the TFR. Galaxies at
absolute magnitudes where such deviations appear are excluded from the
sample.

The initial TFR is calculated using a subset of galaxies dominated by
normal spirals, selected as the core of the distribution via a
two-component GMM fit to the $(x,\,y)$ distribution of galaxies passing
a magnitude pre-filter $-23 < y < -18$ (to remove objects with
implausible absolute magnitudes). The fit accounts for per-galaxy
measurement noise and rectangular survey truncation at the bright and
high-velocity limits of the data. The dominant component represents the
TFR spiral population; the secondary component captures outliers and
dwarf galaxies.

% values in ../output/DR1/selection_ellipse.json
The best-fit GMM yields a dominant-component weight of $w = \GmmW$.
The mean of the dominant component is
$(\bar{x},\,\bar{y}) = (\GmmMeanX,\,\GmmMeanY)$,
with semi-axes $[\GmmSemiA,\,\GmmSemiB]$ and orientation angle $\GmmAngle^\circ$
measured from the $x$-axis (note: the complementary angle appears in plots due to magnitude-axis orientation conventions).
The GMM 1$\sigma$ and 3$\sigma$ ellipses overlaid on the phase-space scatter are shown in
Fig.~\ref{fig:selection_ellipse}.

\begin{figure}
    \centering
    \includegraphics[width=1\linewidth]{\ResultDir/selection_ellipse.png}
    \caption{Distribution of DR2 spiral galaxies used for the initial
    training set selection in $(x,\,y)$ phase space,
    where $x = \log_{10}(V_\text{rot}/100~\text{km\,s}^{-1})$ and $y = M_R$.
    The two GMM ellipses (at 1$\sigma$ and $3\sigma$) are overplotted; the dominant component
    (weight $w = \GmmW$).  The parallelogram defined by the extrema of the
    3$\sigma$ ellipse bound the sample used to determine the initial TFR estimate. }
    \label{fig:selection_ellipse}
\end{figure}

A preliminary selection parallelogram is defined by the $3\sigma$
ellipse of the dominant GMM component.
The slope $\bar{s}$ is fixed by the orientation of the ellipse,
and $\bar{c}_1$, $\bar{c}_2$ are set at $n_{\sigma,\perp} = 3$
perpendicular $\sigma$ from the ellipse center. Values of
$\hat{y}_\text{min}$ and $\hat{y}_\text{max}$ are taken from the
$y$-extent of the $3\sigma$ ellipse. The redshift limits
used in the GMM fit are maintained. The selected parallelogran defining
the initial training region is shown
in Fig.~\ref{fig:selection_ellipse}.

% values in ../output/DR1/select_v2_mle.json
The nominal TFR is given by the maximum likelihood estimator of the
model fit to the data within the parallelogram, yielding a slope of
$\MleSlope$, an intercept of $\MleIntercept$, an intrinsic scatter in $x$
of $\sigma_{\mathrm{int},x} = \MleSigIntX$, and an
intrinsic scatter in $y$ of $\sigma_{\mathrm{int},y} = \MleSigIntY$.

For the full sample of galaxies with $0.03 < z < 0.08$, we predict absolute magnitudes using the
nominal TFR and compute the residual
$\Delta M = \langle M_\text{pred}\rangle - M_\text{obs}$ for every
object. We then calculate the weighted mean residual and its uncertainty
in equal-occupancy bins of $M_R$. Since the nominal TFR slope is only
approximate, we expect a linear trend in this residual profile: an
imperfect slope produces a systematic over- or under-prediction that
varies smoothly with magnitude. Contamination by populations that do not
follow the TFR --- most notably dwarf galaxies --- produces deviations
from this linear trend at the affected magnitude ranges. The magnitude
limits $\hat{y}_\text{min}$ and $\hat{y}_\text{max}$ are set at the
values beyond which the residual profile departs from the linear trend,
identified by visual inspection. The residual profile for the fiducial
selection is shown in Fig.~\ref{fig:pull}; the linear trend is
consistent with the data across the selected range.
The fiducial selection parameters for the DR2 run are listed in
Table~\ref{tab:fiducial} and shown with respect to the full sample in Figure~\ref{fig:data}.


\begin{figure}
    \centering
    \includegraphics[width=1\linewidth]{\ResultDir/select_v2_fiducial_pull.png}
\caption[Weighted-mean residual and pull profile as a function of $M_R$.]{Weighted-mean residual $\Delta M = \langle M_\mathrm{pred}\rangle - M_\mathrm{obs}$
    and pull profile (normalized by bin uncertainty) as a function of $M_R$
    for all catalog objects evaluated at the MLE parameters. Horizontal lines
    show the $\hat{y}_\mathrm{min}$ and $\hat{y}_\mathrm{max}$ range chosen
    to select galaxies with a self-consistent TFR.}
    \label{fig:pull}
\end{figure}

% values in ../output/DR1/select_v2_fiducial.json
\begin{table}
    \centering
    \caption{Fiducial selection parameters for the DR2 training sample.}
    \label{tab:fiducial}
    \begin{tabular}{lc}
        \hline
        Parameter & Value \\
        \hline
        $\hat{y}_\text{min}$ & $\FidHatyMin$ \\
        $\hat{y}_\text{max}$ & $\FidHatyMax$ \\
        $\bar{s}$ & $\FidSlopePlane$ \\
        $\bar{c}_1$ (lower intercept) & $\FidCOne$ \\
        $\bar{c}_2$ (upper intercept) & $\FidCTwo$ \\
        $n_{\sigma,\perp}$ & $\FidNSigPerp$ \\
        $z_\text{min}$ & $\FidZMin$ \\
        $z_\text{max}$ & $\FidZMax$ \\
        Training sample $N$ & $\Ntrain$ \\
        \hline
    \end{tabular}
\end{table}

\begin{figure}
    \centering
    \includegraphics[width=1\linewidth]{\ResultDir/data.png}
    \caption{Full DR2 spiral sample (grey) and training sample (colored points)
    in $(x,\,y)$ phase space.  The oblique selection boundaries and magnitude
    limits are indicated.  The training sample contains $\Ntrain$ galaxies.}
    \label{fig:data}
\end{figure}


\subsection{MCMC Fit}
\label{sec:mcmc}
% fit assessment in ../output/DR1/stansummary.txt, ../output/DR1/diagnoise.txt 
The model is fit to the training sample using Hamiltonian Monte
Carlo as implemented in Stan \citep{STAN_html} via the CmdStan interface.
We run 4 chains, each with 500 warm-up and 1000 sampling iterations.
Convergence diagnostics are satisfactory for all model parameters: $\hat{R} = 1.0$ and
bulk effective sample sizes are reported in Table~\ref{tab:results}.
Stan's \texttt{diagnose} utility confirms no divergent transitions, no treedepth
saturation, and satisfactory E-BFMI across all four chains.
% fit results in  ../output/DR1/stansummary.txt,
The posterior 68\% confidence intervals for the global TFR parameters are
given in Table~\ref{tab:results}.  The full posterior corner plot is shown
in Fig.~\ref{fig:tophat}.

\begin{table}
    \centering
    \caption{Posterior summary for the DR2 tophat model fit.
    Mean and StdDev are computed from the full MCMC chains; 5\%, 50\%, and 95\%
    are posterior quantiles; ESS$_\text{bulk}$ is the bulk effective sample size;
    $\hat{R}$ is the rank-normalized split potential scale reduction factor.}
    \label{tab:results}
    \begin{tabular}{lrrrrrrr}
        \hline
        Parameter & Mean & StdDev & 5\% & 50\% & 95\% & ESS$_\text{bulk}$ & $\hat{R}$ \\
        \hline
        slope $\alpha$         & $\SlopeMean$ & $\SlopeStd$ & $\SlopeQlo$ & $\SlopeQmed$ & $\SlopeQhi$ & $\SlopeEss$ & $\SlopeRhat$ \\
        intercept $\beta$      & $\IntMean$ & $\IntStd$ & $\IntQlo$ & $\IntQmed$ & $\IntQhi$ & $\IntEss$ & $\IntRhat$ \\
        $\sigma_{\mathrm{int},x}$ & $\SigXMean$ & $\SigXStd$ & $\SigXQlo$ & $\SigXQmed$ & $\SigXQhi$ & $\SigXEss$ & $\SigXRhat$ \\
        $\sigma_{\mathrm{int},y}$ & $\SigYMean$ & $\SigYStd$ & $\SigYQlo$ & $\SigYQmed$ & $\SigYQhi$ & $\SigYEss$ & $\SigYRhat$ \\
        \hline
    \end{tabular}
\end{table}

\begin{figure}
    \centering
    \includegraphics[width=1\linewidth]{\ResultDir/tophat.png}
    \caption{Corner plot of the posterior distribution for the four global
    TFR parameters from the DR2 tophat model fit.  Contours show the
    $1\sigma$ and $2\sigma$ confidence regions.}
    \label{fig:tophat}
\end{figure}

\section{Absolute Magnitude Predictions from the TFR}

For all galaxies in the initial DESI DR2 Peculiar Velocity sprial sample,
the TFR determined in
\S\ref{sec:TFRfit} is used to infer their absolute magnitudes. The subset of
galaxies satisfying the magnitude--velocity selection criteria of the training
sample, with no restriction on redshift is referred to as the {\tt MAIN}
sample. Nearly all galaxies in the {\tt MAIN} sample are part of the training set;
nevertheless, we treat each as if it
were out-of-sample. This approximation is justified on two grounds: no single
galaxy has a significant influence on the TFR fit, and accounting for the
covariance between the data and the fitted parameters would substantially
complicate the analysis.



The predicted absolute magnitude for each galaxy is computed via the
posterior predictive distribution derived in
\S\ref{sec:infer_single_galaxy_y}. For a galaxy with measured
rotation-velocity proxy $\hat x_\star$ and uncertainty $\sigma_{x,\star}$,
we adopt a uniform (top-hat) prior on the latent Tully--Fisher magnitude
$y_{\mathrm{TF},\star}$ over the selection interval $[y_{\min}, y_{\max}]$.
Marginalizing over the model parameters $\theta$ via the MCMC posterior
yields the posterior predictive distribution
\begin{equation}
p\!\left(y_\star \mid \hat x_\star, \sigma_{x,\star}, \{\hat x_i, \hat y_i\}\right)
\approx
\frac{1}{M}\sum_{m=1}^{M}
p\!\left(y_\star \mid \hat x_\star, \sigma_{x,\star}, \theta^{(m)}\right),
\label{eq:pred_y_tophat_mcmc}
\end{equation}
where $\{\theta^{(m)}\}_{m=1}^M$ are draws from the MCMC posterior. Each
single-draw predictive is given by the closed form
\begin{equation}
p\!\left(y_\star \mid \hat x_\star, \sigma_{x,\star}, \theta\right)
=
\mathcal{N}\!\left(y_\star \mid \mu_L,\, \sigma_Z^2\right)
\frac{
\Phi\!\left(\frac{y_{\max} - \tilde\mu(y_\star)}{\tilde\sigma}\right)
-
\Phi\!\left(\frac{y_{\min} - \tilde\mu(y_\star)}{\tilde\sigma}\right)
}{
Z_L
},
\label{eq:pred_y_tophat_main}
\end{equation}
where $\Phi(\cdot)$ is the standard normal CDF and $\sigma_{1,\star}^2 = \sigma_{x,\star}^2 + \sigma_{\mathrm{int},x}^2$ as in Eq.~\eqref{eqn:sigma12},
\begin{equation}
\mu_L = \beta + \alpha \hat x_\star,
\qquad
\sigma_Z^2 = \sigma_L^2 + \left(\sigma_{\mathrm{int},y}^{(m)}\right)^2
,
\qquad
\sigma_L = |\alpha|\,\sigma_{1,\star},
\end{equation}
the truncation normalization is
\begin{equation}
Z_L = \Phi\!\left(\frac{y_{\max}-\mu_L}{\sigma_L}\right)
-
\Phi\!\left(\frac{y_{\min}-\mu_L}{\sigma_L}\right),
\end{equation}
and the convolution quantities arising from marginalizing over
$y_{\mathrm{TF},\star}$ are
\begin{equation}
\tilde\sigma^2 = \left(\frac{1}{\sigma_L^2}+\frac{1}{\sigma_{\mathrm{int},y}^2}\right)^{-1},
\qquad
\tilde\mu(y_\star) = \tilde\sigma^2\left(\frac{\mu_L}{\sigma_L^2}
+ \frac{y_\star}{\sigma_{\mathrm{int},y}^2}\right).
\end{equation}




The reported magnitude and its uncertainty for each galaxy are the mean and
standard deviation of the mixture distribution in
Eq.~\eqref{eq:pred_y_tophat_mcmc}. The full posterior predictive covariance
across galaxies is obtained via the law of total covariance. Conditional on
the TFR parameters $\theta$, galaxy predictions are independent, so
off-diagonal covariance arises solely from the shared uncertainty in $\theta$.
For each MCMC draw $\theta^{(m)}$, define the per-draw conditional mean
\begin{equation}
    \mu_i^{(m)} \;=\; \mathbb{E}\left[y_i \mid \hat{x}_i, \sigma_{1,i}^{(m)},
    \theta^{(m)}\right],
\end{equation}
which is the truncated-normal mean of $y_{\mathrm{TF},i}$ given $\hat{x}_i$
and $\theta^{(m)}$ (Eq.~\eqref{eq:post_yTF_tophat}), and the per-draw
predictive variance
\begin{equation}
    v_i^{(m)} \;=\; \operatorname{Var}\!\left(y_{\mathrm{TF},i} \mid \hat{x}_i,
    \theta^{(m)}\right) \;+\; \left(\sigma_{\mathrm{int},y}^{(m)}\right)^2,
\end{equation}
including both flavors of intrinsic scatter: $\sigma_{\mathrm{int},y}$ enters
additively, while $\sigma_{\mathrm{int},x}$ enters the truncated-normal width
through $\sigma_L = |\alpha|\,\sigma_{1,\star}$, where
$\sigma_{1,\star}^2 = \sigma_{x,\star}^2 + \sigma_{\mathrm{int},x}^2$.
The diagonal of the predictive covariance then decomposes into three
contributions: calibration variance ($\operatorname{Cov}_\theta$ of the
conditional means), expected predictive variance
($\mathbb{E}_\theta[v_i^{(m)}]$, encompassing intrinsic scatter and
truncated-normal variance), and measurement noise $\sigma_{y,i}^2$ in
$\hat{y}_i$, where $\bar{y}_i = \frac{1}{M}\sum_m \mu_i^{(m)}$ is the
posterior predictive mean. Off-diagonal elements follow from the
calibration-variance term alone, exploiting conditional independence.
The full posterior predictive covariance matrix for a subsample of
galaxies ordered by redshift is visualized in Fig.~\ref{fig:tophat_cov}.  


\begin{figure}
    \centering
    \includegraphics[width=1\linewidth]{\ResultDir/tophat_cov_sub_noobs.png}
    \caption{Visualization of the absolute magnitude posterior predictive covariance
    matrix
    for a subsample of galaxies ordered by redshift.  Off-diagonal
    correlations arise from the shared TFR calibration parameters.
    The diagonal for the matrix of the {\tt MAIN} sample is not clearly visible at the resolution of this plot.}
    \label{fig:tophat_cov}
\end{figure}

\section{Peculiar Velocities/Magnitudes}
\label{sec:peculiar_magnitudes}

Peculiar velocities are manifest as peculiar magnitudes, i.e.\ deviations
between observed and TFR-predicted absolute magnitudes.
These residuals $\hat{y}_\text{pred} - \hat{y}_\text{obs}$
are calculated for the DR2 PV spiral sample.

We define a grid in $(x,\,y)$ phase space and compute the mean
residual for each element.
Figure~\ref{fig:tophat_grid} shows this residual grid, both for the full sample
and for the {\tt MAIN} sample.

As discussed in \S\ref{sec:training_selection}, the magnitude limits
$\hat{y}_\text{min}$ and $\hat{y}_\text{max}$ were chosen to exclude
regions where the TFR fit is not self-consistent, by examining the combined magnitude
residuals in slices of redshift.
The left panel of Fig.~\ref{fig:tophat_grid} illustrates why this is
necessary: for galaxies brighter than $M_R = -19$, the zero-residual
locus follows the TFR slope, with residuals increasing monotonically
with distance from it, consistent with Gaussian scatter and measurement
uncertainties. For galaxies fainter than $M_R = -19$, however, the
zero-residual locus deviates from this slope, indicating a population
--- likely dwarf galaxies --- that does not follow the TFR.
A similar but more sublte 
deviation is observed at the bright end, which could be due to galaxy outliers
or limitations of the top-hat model in this regime.
The right panel of Fig.~\ref{fig:tophat_grid} shows that for
the $\tt MAIN$ sample, for which selection criteria are applied (but no redshift restriction),
the residuals are consistent with zero across the full phase space, with no evidence for a
non-linear TFR or magnitude-dependent bias.
There is a slight asymmetry in the residuals, as the sample selection parallelogram
is not perfectly centered or aligned with the TFR fit to the training sample and used
in the selection of the final sample.

\begin{figure}
    \centering
    \includegraphics[width=0.49\linewidth]{\ResultDir/tophat_grid_full.png}
    \includegraphics[width=0.49\linewidth]{\ResultDir/tophat_grid.png}
    \caption{Mean prediction residual $\hat{y}_\text{pred} - \hat{y}_\text{obs}$
    binned on a grid in $(x,\,y)$ phase space for the training sample.  Left:
    the DR2 PV spiral, Right: the {\tt MAIN} sample.
    The residuals are consistent with zero, with no evidence for systematic
    trends.}
    \label{fig:tophat_grid}
\end{figure}

The importance of redshift is demonstrated in Fig.~\ref{fig:redshift_tophat}, which shows\
the residuals as a function of redshift for both the full sample and the {\tt MAIN} sample. 
For the full sample, there is a clear redshift dependence in the residuals, with a systematic
trend that deviates from zero at low and high redshifts.
This indicates that the TFR fit is not self-consistent across the full redshift range,
likely due to contamination by non-TFR-following populations (e.g., dwarf galaxies) at certain redshifts.
In contrast, for the {\tt MAIN} sample, which is selected based on the criteria defined
in \S\ref{sec:training_selection}, the residuals are consistent with zero across the entire redshift range,
at least to first-order, demonstrating that the selection criteria mitigate
redshift-dependent systematic effects.


\begin{figure}
    \centering
    \includegraphics[width=1\linewidth]{\ResultDir/redshift_tophat.png}
    \caption{TFR residuals as a function of observed redshift for the
     input and {\tt MAIN} samples.}
    \label{fig:redshift_tophat}
\end{figure}







\subsection{Data Products}
The catalog is provided as \texttt{tophat\_catalog.fits}, which augments the
input FITS table with five new columns:
\begin{itemize}
    \item \texttt{MU\_TF}: the TFR-inferred distance modulus
    $\mu_{TF} = m_{R,\text{corr}} - \langle M_*\rangle_\text{pred}$,
    where $m_{R,\text{corr}}$ is the corrected apparent $R$-band magnitude
    and $\langle M_*\rangle_\text{pred}$ is the posterior predictive mean
    absolute magnitude from the TFR model (float32).
    \item \texttt{MU\_ERR}: the combined uncertainty
    $\sigma_{\mu_{TF}} = \sqrt{\sigma_{m,\text{corr}}^2 + \sigma_\text{pred}^2}$,
    where $\sigma_{m,\text{corr}}$ is the observational magnitude uncertainty
    and $\sigma_\text{pred}$ is the posterior predictive standard deviation (float32).
    \item \texttt{LOGDIST}: the log-distance ratio
    $\eta = 0.2\,(\mu_z - \mu_{TF})$,
    where $\mu_z = m_R - M_R(z)$ is the CMB-frame distance modulus
    derived from the observed redshift.
    Equivalently, $\eta = \log_{10}(d_H/d_{TF})$, the base-10 logarithm of
    the ratio of the Hubble distance to the TFR distance (float32).
    \item \texttt{LOGDIST\_ERR}: the uncertainty on \texttt{LOGDIST},
    $\sigma_\eta = 0.2\,\sigma_{\mu_{TF}}$ (float32).
    \item \texttt{MAIN}: boolean flag (\texttt{True}/\texttt{False}) indicating
    whether the galaxy satisfies the phase-space selection cuts
    (absolute magnitude range and oblique TFR-plane boundaries defined in
    \texttt{config.json}), evaluated without a redshift restriction.
    Only \texttt{MAIN}$=$\texttt{True} galaxies are included in
    \texttt{tophat\_cov.fits}.
\end{itemize}

The posterior predictive covariance matrix is provided as
\texttt{tophat\_cov.fits}.  The file contains a single $N\times N$ float32
image extension, where $N$ is the number of \texttt{MAIN}-flagged galaxies.
Rows and columns are ordered to match the \texttt{MAIN}$=$\texttt{True} rows
of \texttt{tophat\_catalog.fits} in their original catalog order.  Entry
$(g,g')$ gives $\mathrm{Cov}(\mu_{TF}^{(g)},\,\mu_{TF}^{(g')})$, the
posterior predictive covariance of the TF distance moduli; diagonal entries
additionally include the per-galaxy observational magnitude variance
$\sigma_{m,\text{obs}}^2$.  This matrix encodes the correlated uncertainties
arising from the shared TFR calibration parameters and is required for
cosmological analyses that propagate correlated peculiar velocity errors.
The matrix structure is visualized in Fig.~\ref{fig:tophat_cov}.




\section{Conclusions}
\label{sec:conclusions}
We have presented a model-based calibration of the Tully-Fisher relation
using the full spiral-galaxy sample from the DESI DR2 Peculiar Velocity survey.
The key results are as follows.
\begin{itemize}
    \item A two-component GMM fit to the DR2 phase space identifies the
    dominant spiral TFR population (weight $w = \GmmW$) and defines the
    orientation of oblique selection cuts.  The fiducial selection yields
    a training sample of $\Ntrain$ galaxies with a flat pull profile confirming
    no selection-induced bias.
    \item The Stan tophat model, which places a uniform prior on the latent
    log-velocity within the observed selection region, correctly accounts for
    the complex oblique selection boundary without the need for explicit
    bias corrections.
    \item MCMC posterior parameters are slope $\alpha = \SlopeMean \pm \SlopeStd$,
    intercept $\beta = \IntMean \pm \IntStd$, and intrinsic scatter
    $\sigma_{\mathrm{int},x} = \SigXMean \pm \SigXStd$,
    $\sigma_{\mathrm{int},y} = \SigYMean \pm \SigYStd$.
    Residual diagnostics in both phase space and redshift show no evidence
    for systematic trends.
    \item The calibrated model is applied to produce a peculiar velocity
    catalog (\texttt{tophat\_catalog.fits}) and a posterior predictive
    covariance matrix (\texttt{tophat\_cov.fits}) for the full MAIN sample.
\end{itemize}

This work is distinguished from other approaches in that it fits a probabilistic
TFR model rather than fitting a line directly to the data.  The key distinction
is that parameter estimation is based on the model likelihood, whereas other
approaches optimize a reasonable but essentially arbitrary merit function.
A feature of the TFR model is that it must specify the latent parameter
distribution, and we find that the results are sensitive to this choice.
An advantage of our approach is that the model can be applied to independent
datasets in a self-consistent manner.

Peculiar velocity measurements from the MAIN catalog, including the study
of large-scale flows and comparison with existing peculiar velocity compilations,
will be presented in a follow-up paper.


\begin{acknowledgments}
    This material is based upon work supported by the U.S. Department of Energy (DOE), Office of Science, Office of High-Energy Physics, under Contract No. DE--AC02--05CH11231, and by the National Energy Research Scientific Computing Center, a DOE Office of Science User Facility under the same contract. Additional support for DESI was provided by the U.S.
    National Science Foundation (NSF), Division of Astronomical Sciences under Contract No. AST-0950945 to the NSF's National Optical-Infrared Astronomy Research Laboratory; the Science and Technology Facilities Council of the United Kingdom; the Gordon and Betty Moore Foundation;
    the Heising-Simons Foundation; the French Alternative Energies and Atomic Energy Commission (CEA); the National Council of Humanities, Science and Technology of Mexico (CONAHCYT); the Ministry of Science, Innovation and Universities of Spain
    (MICIU/AEI/10.13039/501100011033), and by the DESI Member Institutions: \url{https://www.desi.lbl.gov/collaborating-institutions}. Any opinions, findings, and conclusions or recommendations expressed in this material are those of the author(s) and do not necessarily reflect the views of the U. S. National Science Foundation, the U. S. Department of Energy, or any of the listed funding agencies.

The authors are honored to be permitted to conduct scientific research on Iolkam Du'ag (Kitt Peak), a mountain with particular significance to the Tohono O'odham Nation.
\end{acknowledgments}

% Software: Stan \citep{STAN_html} using the CmdStan interface,
% ChainConsumer \citep{Hinton2016}


\appendix
\input{../doc/model1}

\input{../doc/model2}


\section{Color-Correction Model}
\label{sec:color_correction}

\subsection{Motivation}

Empirical diagnostics of the TFR fit reveal color-dependent residuals: galaxies
with redder $r-z$ colors are systematically brighter than predicted by the
baseline model at fixed rotation velocity.
We extend the model by introducing a single global parameter $\gamma$, the slope
of the luminosity--color relation at fixed rotation velocity, so that a galaxy
with $r-z$ color $c_i$ redder than the population mean $\mu_c$ has
$r$-band absolute magnitude shifted by $\gamma\eta_i$, where $\eta_i$ is the
component of color that is independent of velocity.

\subsection{Generative Model}
\label{sec:cc:generative}

\paragraph{Latent variables.}
Each galaxy $i$ has latent variables $y_{\mathrm{TF},i}$, $x_i$, and $c_i$
(intrinsic color).
The generative model is
\begin{align}
y_{\mathrm{TF},i} &\sim \mathrm{Uniform}(y_{\min},\,y_{\max}), \label{eq:cc:yTF_prior}\\
x_i \mid y_{\mathrm{TF},i} &\sim \mathcal{N}\!\left(\frac{y_{\mathrm{TF},i}-\beta}{\alpha},\;\sigma_{\mathrm{int},x}^2\right), \label{eq:cc:x_given_yTF}\\
c_i \mid x_i &\sim \mathcal{N}\!\left(\mu_c + \delta(x_i - \bar{x}),\;\tau_c^2\right), \label{eq:cc:c_given_x}\\
\eta_i &\equiv c_i - \mu_c - \delta(x_i - \bar{x}), \label{eq:cc:eta_def}\\
y_i \mid y_{\mathrm{TF},i},\,\eta_i &\sim \mathcal{N}\!\left(y_{\mathrm{TF},i} + \gamma\eta_i + \Delta_{r,i},\;\sigma_{\mathrm{int},y}^2\right), \label{eq:cc:y_given_yTF_eta}\\
z_i \mid y_{\mathrm{TF},i},\,\eta_i,\,x_i &\sim \mathcal{N}\!\left(y_{\mathrm{TF},i} + (\gamma-1)\eta_i - \mu_c - \delta(x_i-\bar{x}) + \Delta_{z,i},\;\sigma_{\mathrm{int},z}^2\right), \label{eq:cc:z_given_yTF_eta}
\end{align}
with measurement equations
\begin{align}
\hat{x}_i \mid x_i &\sim \mathcal{N}(x_i,\;\sigma_{x,i}^2), \label{eq:cc:xhat}\\
\hat{y}_i \mid y_i &\sim \mathcal{N}(y_i,\;\sigma_{y,i}^2), \label{eq:cc:yhat}\\
\hat{z}_i \mid z_i &\sim \mathcal{N}(z_i,\;\sigma_{z,i}^2). \label{eq:cc:zhat}
\end{align}
Here $\bar{x}$ is the sample mean of $\hat{x}$ and
\begin{align}
\Delta_{r,i} &\equiv \alpha_{k,r} \bigl[\ln(1 + z_{\mathrm{obs},i}) - \overline{\ln(1+z)}\bigr], \label{eq:cc:Delta_r}\\
\Delta_{z,i} &\equiv \alpha_{k,z} \bigl[\ln(1 + z_{\mathrm{obs},i}) - \overline{\ln(1+z)}\bigr]  \label{eq:cc:Delta_z}
\end{align}
are band-dependent residual $k$-corrections that absorb any
redshift-dependent offsets not already removed by the upstream photometric
corrections, where
$\overline{\ln(1+z)} \equiv N^{-1}\sum_j \ln(1+z_{\mathrm{obs},j})$ is the
sample mean.  Centering the regressors removes degeneracies with the
intercept and $\mu_c$; the mean $k$-corrections are absorbed by these
parameters by construction.
Separate slopes are needed because the $r$- and $z$-band $k$-corrections
differ: the observed $r{-}z$ color reddens by $\sim\!0.12$~mag from
$z_{\mathrm{obs}}=0.02$ to $0.1$ in the training sample, a trend that a
single shared $\Delta_i$ cannot capture.
The variable $\eta_i \sim \mathcal{N}(0,\tau_c^2)$ is the residual color after
removing the population-level color--velocity slope $\delta$; it is
independent of $x_i$ by construction.

\paragraph{Physical interpretation.}
The parameter $\gamma$ is the slope of the luminosity--color relation at fixed
rotation velocity: at fixed $x_i$, a galaxy one unit redder in residual color
($\eta_i$ increased by 1) is $\gamma$ magnitudes brighter ($y_i$ more negative
by $|\gamma|$ if $\gamma < 0$).
The parameter $\delta$ captures the population-level trend that faster rotators
tend to be redder; it is not a distance-indicator correction but rather a
description of the galaxy population.

\paragraph{Identification.}
The model is identified because $\gamma$ enters only through the residual color
$\eta_i$, which is orthogonal to $x_i$ by construction (Eq.~\ref{eq:cc:eta_def}).
Specifically:
\begin{enumerate}
\item The transformation $c_i \to c_i + e$, $\beta \to \beta - \gamma e$ changes
  the prior mean of $c_i$ from $\mu_c$ to $\mu_c + e$, which is penalized by the
  prior $c_i \mid x_i \sim \mathcal{N}(\mu_c + \delta(x_i-\bar{x}),\tau_c^2)$
  unless $\mu_c$ is simultaneously shifted.
  The prior $\mu_c \sim \mathcal{N}(\bar{c}_{\mathrm{obs}}, 1)$ (below) breaks
  this near-degeneracy.
\item The transformation $c_i \to c_i + d\,x_i$, $\alpha \to \alpha - \gamma d$
  would require the color--velocity slope to shift by $d$, but $\gamma$ acts on
  $\eta_i$ (residual color at fixed $x_i$), not on $c_i$ directly; the linear-in-$x$
  component of color is absorbed by $\delta$, not by $\alpha$.
  Because $\eta_i$ is defined using the latent $x_i$ rather than $\hat{x}_i$,
  no attenuation bias arises.
\end{enumerate}

\paragraph{Priors.}
The parameters $\alpha$, $\beta$, $\sigma_{\mathrm{int},x}$, $\sigma_{\mathrm{int},y}$
inherit the priors of the baseline model (Appendix~\ref{sec:likelihood}).
Additional priors:
\begin{align}
\gamma &\sim \mathcal{N}(0,\,1), \label{eq:cc:gamma_prior}\\
\delta &\sim \mathcal{N}(0,\,1), \label{eq:cc:delta_prior}\\
\mu_c  &\sim \mathcal{N}(\bar{c}_{\mathrm{obs}},\,1), \label{eq:cc:muc_prior}\\
\tau_c &\sim \mathrm{Half\text{-}Cauchy}(0,\,0.3)\;\text{truncated to }\tau_c > 0.014, \label{eq:cc:tauc_prior}\\
\ln\sigma_{\mathrm{int},z} &\sim \mathcal{N}(-3,\,2), \label{eq:cc:sigz_prior}\\
\alpha_{k,r} &\sim \mathcal{N}(0,\,5), \label{eq:cc:alphakr_prior}\\
\alpha_{k,z} &\sim \mathcal{N}(0,\,5), \label{eq:cc:alphakz_prior}
\end{align}
where $\bar{c}_{\mathrm{obs}} \equiv N^{-1}\sum_i (\hat{y}_i - \hat{z}_i)$ is
the sample mean observed color.
The prior on $\gamma$ is not truncated in the implementation.
Without a hard constraint, the likelihood admits a near-degenerate solution
at $\gamma > 0$ in which the roles of $\sigma_{\mathrm{int},y}$ and
$\sigma_{\mathrm{int},z}$ are exchanged; fits to the data show that
the $\gamma < 0$ mode has substantially higher posterior probability, and in
practice the sampler concentrates on $\gamma < 0$ without needing the bound
enforced explicitly.
The lower bound $\tau_c > 0.014$ excludes a second degenerate mode at
$\tau_c \approx 0.006$--$0.013$ in which the color correction effectively
shuts off ($\gamma\tau_c \to 0$); this mode occupies a clearly separated
region of the posterior with substantially lower probability than the
physical mode at $\tau_c \approx 0.03$.

\paragraph{Sampling parameterization.}
Hamiltonian Monte Carlo on the natural $(\gamma,\tau_c)$ coordinates is
hampered by two coupled features of the posterior geometry.  First, the
latent contribution to the $r$-band magnitude,
$\gamma\eta_i \sim \mathcal{N}(0,p^2)$ with $p \equiv \gamma\tau_c$,
depends on $\gamma$ and $\tau_c$ only through their product.  The $y$-data
alone are therefore degenerate along contours of constant $p$; the $z$-
and color-data resolve $\gamma$ and $\tau_c$ separately only weakly,
leaving a banana-shaped posterior ridge.  Second, the prior on $\tau_c$
places appreciable mass near zero, and the joint posterior contracts
sharply as $\tau_c \to 0$: a single HMC step size cannot simultaneously
resolve the narrow region near $\tau_c = 0$ and the broader region away
from it.

We mitigate both pathologies by sampling
\begin{equation}
(p,\,\ell)\;\equiv\;(\gamma\tau_c,\;\log\tau_c)
\label{eq:cc:reparam}
\end{equation}
in place of $(\gamma,\tau_c)$, with $\gamma = p\,e^{-\ell}$ and
$\tau_c = e^{\ell}$ recovered as derived quantities.  The $p$-coordinate
aligns one sampler axis with the banana ridge; the log scale on $\tau_c$
renders the geometry scale-invariant near $\tau_c = 0$, removing the
contraction.  The change of variables has unit Jacobian determinant,
\begin{equation}
\left|\det\frac{\partial(p,\ell)}{\partial(\gamma,\tau_c)}\right|
= \left|\det\begin{pmatrix}\tau_c & \gamma\\[2pt] 0 & \tau_c^{-1}\end{pmatrix}\right|
= 1,
\label{eq:cc:reparam_jac}
\end{equation}
because the $\tau_c^{-1}$ Jacobian from $\gamma\to p$ at fixed $\tau_c$ is
exactly cancelled by the $\tau_c$ Jacobian from $\tau_c\to\ell$.  The
priors of Eqs.~\eqref{eq:cc:gamma_prior}--\eqref{eq:cc:alphakz_prior}
therefore apply on the derived $\gamma$ and $\tau_c$ with no additional
volume correction.

\subsection{Marginalization over Latent Variables}
\label{sec:cc:marginalization}

\paragraph{Marginalizing $\eta_i$.}
For fixed $(y_{\mathrm{TF},i}, x_i)$, the residual color
$\eta_i \sim \mathcal{N}(0,\tau_c^2)$ is independent of $x_i$.
Since $y_i = y_{\mathrm{TF},i} + \gamma\eta_i + \Delta_{r,i} + \epsilon_y$ (with $\epsilon_y \sim \mathcal{N}(0,\sigma_{\mathrm{int},y}^2)$)
and $z_i = y_{\mathrm{TF},i} + (\gamma-1)\eta_i - \mu_c - \delta(x_i-\bar{x}) + \Delta_{z,i} + \epsilon_z$
(with $\epsilon_z \sim \mathcal{N}(0,\sigma_{\mathrm{int},z}^2)$ independent of $\epsilon_y$), we have
\begin{align}
\tilde{y}_i &\equiv y_i - y_{\mathrm{TF},i} - \Delta_{r,i} = \gamma\eta_i + \epsilon_y,\\
\tilde{z}_i &\equiv z_i - [y_{\mathrm{TF},i} + \Delta_{z,i} - \mu_c - \delta(x_i-\bar{x})] = (\gamma-1)\eta_i + \epsilon_z.
\end{align}
Marginalizing $\eta_i \sim \mathcal{N}(0,\tau_c^2)$ and adding measurement noise
gives the joint marginal of $(\hat{y}_i, \hat{z}_i)$
given $(y_{\mathrm{TF},i}, x_i)$:
\begin{equation}
\begin{pmatrix}\hat{y}_i \\ \hat{z}_i\end{pmatrix}
\;\Bigg|\; y_{\mathrm{TF},i},\,x_i
\;\sim\;
\mathcal{N}\!\left(
\begin{pmatrix}
  y_{\mathrm{TF},i} + \Delta_{r,i}\\[2pt]
  y_{\mathrm{TF},i} + \Delta_{z,i} - \mu_c - \delta(x_i-\bar{x})
\end{pmatrix},\;
\boldsymbol{A}_i
\right),
\label{eq:cc:yz_given_yTF_x}
\end{equation}
where the covariance matrix is
\begin{equation}
\boldsymbol{A}_i =
\begin{pmatrix}
  \gamma^2\tau_c^2 + \sigma_{\mathrm{int},y}^2 + \sigma_{y,i}^2 &
  \gamma(\gamma-1)\tau_c^2 \\[4pt]
  \gamma(\gamma-1)\tau_c^2 &
  (\gamma-1)^2\tau_c^2 + \sigma_{\mathrm{int},z}^2 + \sigma_{z,i}^2
\end{pmatrix}.
\label{eq:cc:A}
\end{equation}
The off-diagonal $\gamma(\gamma-1)\tau_c^2$ arises
because $\hat{y}$ and $\hat{z}$ share the color residual $\eta_i$
(with coefficients $\gamma$ and $\gamma-1$); the intrinsic scatters
$\epsilon_y$ and $\epsilon_z$ are independent and do not contribute.

\paragraph{Sanity check: $\gamma = 0$, $\alpha_{k,r} = \alpha_{k,z} = 0$.}
Setting $\gamma = 0$ and $\alpha_{k,r} = \alpha_{k,z} = 0$: $\boldsymbol{A}_i = \mathrm{diag}(\sigma_{\mathrm{int},y}^2+\sigma_{y,i}^2,\; \sigma_{\mathrm{int},z}^2 + \sigma_{z,i}^2)$
and the off-diagonal vanishes, i.e., $\hat{y}$ and
$\hat{z}$ are independent when color has no
effect on $r$-band luminosity.
The marginal for $\hat{y}_i$ reduces to
$\mathcal{N}(y_{\mathrm{TF},i}, \sigma_{\mathrm{int},y}^2 + \sigma_{y,i}^2)$,
recovering the baseline model of Appendix~\ref{sec:likelihood}.

\paragraph{Marginalizing $x_i$.}
Adding the measurement equation $\hat{x}_i \mid x_i \sim \mathcal{N}(x_i, \sigma_{x,i}^2)$
and the model $x_i \mid y_{\mathrm{TF},i} \sim \mathcal{N}((y_{\mathrm{TF},i}-\beta)/\alpha,\,\sigma_{\mathrm{int},x}^2)$,
the variable $x_i$ enters the mean of $\hat{z}_i$ linearly through the
$\delta(x_i-\bar{x})$ term. Marginalizing $x_i$ produces a trivariate Gaussian
for $(\hat{x}_i, \hat{y}_i, \hat{z}_i)$ given $y_{\mathrm{TF},i}$.
Defining
\begin{align}
\sigma_{1,i}^2 &\equiv \sigma_{x,i}^2 + \sigma_{\mathrm{int},x}^2,\\
\mu_{x,i}(y_{\mathrm{TF}}) &\equiv \frac{y_{\mathrm{TF}}-\beta}{\alpha},
\end{align}
the trivariate marginal is
\begin{equation}
\begin{pmatrix}\hat{x}_i \\ \hat{y}_i \\ \hat{z}_i\end{pmatrix}
\;\Bigg|\; y_{\mathrm{TF},i}
\;\sim\;
\mathcal{N}\!\left(
\begin{pmatrix}
  \mu_{x,i}(y_{\mathrm{TF},i})\\[2pt]
  y_{\mathrm{TF},i} + \Delta_{r,i}\\[2pt]
  y_{\mathrm{TF},i} + \Delta_{z,i} - \mu_c - \delta(\mu_{x,i}(y_{\mathrm{TF},i})-\bar{x})
\end{pmatrix},\;
\boldsymbol{B}_i
\right),
\label{eq:cc:trivariate}
\end{equation}
where the $3\times 3$ covariance matrix is
\begin{equation}
\boldsymbol{B}_i =
\begin{pmatrix}
  \sigma_{1,i}^2 & 0 & -\delta\sigma_{\mathrm{int},x}^2 \\[4pt]
  0 & A_{11} & A_{12} \\[4pt]
  -\delta\sigma_{\mathrm{int},x}^2 & A_{12} & A_{22} + \delta^2\sigma_{\mathrm{int},x}^2
\end{pmatrix},
\label{eq:cc:B}
\end{equation}
with $A_{11}$, $A_{12}$, $A_{22}$ the entries of $\boldsymbol{A}_i$
(Eq.~\ref{eq:cc:A}).
The $(\hat{x}, \hat{y})$ off-diagonal is zero because $x_i$ does not enter $y_i$
once $y_{\mathrm{TF},i}$ is fixed; the $(\hat{x}, \hat{z})$ off-diagonal is
$-\delta\sigma_{\mathrm{int},x}^2$ because $\hat{z}_i$ depends on the latent $x_i$
(not $\hat{x}_i$), and only the intrinsic variance $\sigma_{\mathrm{int},x}^2$ is
shared between $\hat{x}_i$ and $\hat{z}_i$.

The remaining 1D integral over $y_{\mathrm{TF},i}$ admits a closed form because the
mean vector in Eq.~\eqref{eq:cc:trivariate} is linear in $y_{\mathrm{TF},i}$. Write
\begin{equation}
\boldsymbol{\mu}(y_{\mathrm{TF}}) = \boldsymbol{a}_i + \boldsymbol{b}\,y_{\mathrm{TF}},
\quad
\boldsymbol{a}_i =
\begin{pmatrix}
  -\beta/\alpha \\
  \Delta_{r,i} \\
  \Delta_{z,i} - \mu_c + \delta\bar{x} + \delta\beta/\alpha
\end{pmatrix},
\quad
\boldsymbol{b} =
\begin{pmatrix}
  1/\alpha \\
  1 \\
  1 - \delta/\alpha
\end{pmatrix}.
\label{eq:cc:linear_mean}
\end{equation}
Let $\hat{\boldsymbol{d}}_i \equiv (\hat{x}_i,\hat{y}_i,\hat{z}_i)^T - \boldsymbol{a}_i$ and define the scalars
\begin{equation}
\xi_i \equiv \boldsymbol{b}^T \boldsymbol{B}_i^{-1} \boldsymbol{b},
\quad
\phi_i \equiv \boldsymbol{b}^T \boldsymbol{B}_i^{-1} \hat{\boldsymbol{d}}_i,
\quad
\chi_i \equiv \hat{\boldsymbol{d}}_i^T \boldsymbol{B}_i^{-1} \hat{\boldsymbol{d}}_i.
\label{eq:cc:scalars}
\end{equation}
Completing the square in $y_{\mathrm{TF}}$,
\begin{equation}
(\hat{\boldsymbol{d}}_i - \boldsymbol{b}\,y_{\mathrm{TF}})^T \boldsymbol{B}_i^{-1}
(\hat{\boldsymbol{d}}_i - \boldsymbol{b}\,y_{\mathrm{TF}})
= \xi_i\,(y_{\mathrm{TF}} - \mu^*_i)^2 + Q_{0,i},
\label{eq:cc:complete_square}
\end{equation}
with $\mu^*_i \equiv \phi_i/\xi_i$ and $Q_{0,i} \equiv \chi_i - \phi_i^2/\xi_i$.
The integrand becomes a 1D Gaussian in $y_{\mathrm{TF}}$ times a $y_{\mathrm{TF}}$-independent prefactor; integrating over $[y_{\min}, y_{\max}]$ gives
\begin{equation}
\int_{y_{\min}}^{y_{\max}}
\mathcal{N}_3\!\bigl((\hat{x}_i,\hat{y}_i,\hat{z}_i)^T;\,\boldsymbol{\mu}(y_{\mathrm{TF}}),\,\boldsymbol{B}_i\bigr)\,
dy_{\mathrm{TF}}
=
\frac{e^{-Q_{0,i}/2}}{2\pi\sqrt{|\boldsymbol{B}_i|\,\xi_i}}\,
\bigl[\Phi(u^{\max}_i) - \Phi(u^{\min}_i)\bigr],
\label{eq:cc:numerator_closed}
\end{equation}
where $u^{\bullet}_i \equiv \sqrt{\xi_i}\,(y_{\bullet} - \mu^*_i)$ and $\Phi$ is the
standard normal CDF. Equation~\eqref{eq:cc:numerator_closed} supersedes the
sinh-reparameterization quadrature of Appendix~\ref{sec:likelihood} for the
trivariate likelihood numerator; the selection integral in
Section~\ref{sec:cc:selection} still requires 1D quadrature, since the
denominator includes an additional bivariate integration over the selection
region.

\subsection{Selection Probability}
\label{sec:cc:selection}

The selection function $S_i = 1$ depends on the observed $(\hat{x}_i, \hat{y}_i)$
only; the $z$-band observation $\hat{z}_i$ is not used in the selection cuts
(Eq.~10 of the main text).
The selection probability is therefore
\begin{equation}
P(S_i=1 \mid y_{\mathrm{TF},i},\,\theta,\,\gamma)
= \int\!\!\!\int_{(\hat{x},\hat{y})\in\mathcal{S}}
  p(\hat{x}_i,\hat{y}_i \mid y_{\mathrm{TF},i},\,\theta,\,\gamma)\,
  d\hat{x}_i\,d\hat{y}_i,
\label{eq:cc:Psel}
\end{equation}
where $p(\hat{x}_i,\hat{y}_i \mid y_{\mathrm{TF},i}, \theta)$ is the
$(\hat{x},\hat{y})$ marginal of Eq.~\eqref{eq:cc:trivariate}, i.e.\ a bivariate
Gaussian with mean $(\mu_{x,i}(y_{\mathrm{TF},i}),\; y_{\mathrm{TF},i} + \Delta_{r,i})^T$ and
covariance $\bigl(\begin{smallmatrix}\sigma_{1,i}^2 & 0\\ 0 & A_{11}\end{smallmatrix}\bigr)$.
The $\hat{z}$ component does not enter because it is independent of $\hat{x}$
and $\hat{y}$ given $y_{\mathrm{TF},i}$; the mean shift from $\delta$ is
carried only by $\hat{z}$, not by $\hat{y}$.
Since $\Delta_{r,i}$ shifts $\mathrm{E}[\hat{y}_i \mid y_{\mathrm{TF},i}]$ by a
galaxy-dependent amount, it enters the selection probability: at fixed
$y_{\mathrm{TF},i}$, a non-zero $\Delta_{r,i}$ moves the predicted $\hat{y}_i$
relative to the fixed boundaries $\hat{y}_{\min}$, $\hat{y}_{\max}$.
Equivalently, one may evaluate the selection integral with $\Delta_{r,i}=0$ but
with shifted boundaries $\hat{y}_{\min} - \Delta_{r,i}$ and
$\hat{y}_{\max} - \Delta_{r,i}$.
Substituting into the sheared-coordinate bivariate CDF machinery of
Appendix~\ref{sec:likelihood} (Eq.~\ref{eq:Ii_def}), with
$\sigma_{2,i}^2$ replaced by $A_{11} = \gamma^2\tau_c^2 + \sigma_{\mathrm{int},y}^2 + \sigma_{y,i}^2$,
gives the selection probability as a 1D integral over $y_{\mathrm{TF},i}$.

\subsection{Posterior Predictive Distribution}
\label{sec:cc:predict}

For a new galaxy with measurements $(\hat{x}_\star, \hat{z}_\star)$ and uncertainties
$(\sigma_{x,\star}, \sigma_{z,\star})$, the posterior predictive for $\hat{y}_\star$ at
fixed $(\theta, \gamma)$ is obtained by marginalizing $y_{\mathrm{TF},\star}$,
$x_\star$, and $\eta_\star$.

\paragraph{Marginalizing $\eta_\star$ and $x_\star$.}
Given $y_{\mathrm{TF},\star}$, the trivariate marginal
Eq.~\eqref{eq:cc:trivariate} gives the joint distribution of
$(\hat{x}_\star, \hat{y}_\star, \hat{z}_\star)$.
Conditioning on the observed $(\hat{x}_\star, \hat{z}_\star)$ via the
Schur complement gives
\begin{equation}
\hat{y}_\star \mid \hat{x}_\star,\,\hat{z}_\star,\,y_{\mathrm{TF},\star}
\;\sim\;
\mathcal{N}\!\bigl(\mu_{y|\hat{x}\hat{z}}(y_{\mathrm{TF},\star}),\;\sigma_{y|\hat{x}\hat{z}}^2\bigr),
\label{eq:cc:y_given_xz_yTF}
\end{equation}
where
\begin{align}
\mu_{y|\hat{x}\hat{z}}(y_{\mathrm{TF},\star})
&= y_{\mathrm{TF},\star} + \Delta_{r,\star}
  + \boldsymbol{b}_{\mathrm{reg}}^T
    \begin{pmatrix}
      \hat{x}_\star - \mu_{x,\star}(y_{\mathrm{TF},\star})\\[2pt]
      \hat{z}_\star - [y_{\mathrm{TF},\star} + \Delta_{z,\star} - \mu_c - \delta(\mu_{x,\star}(y_{\mathrm{TF},\star})-\bar{x})]
    \end{pmatrix},
\label{eq:cc:mu_ygivenxz}\\
\sigma_{y|\hat{x}\hat{z}}^2
&= A_{11} - \boldsymbol{b}_{\mathrm{reg}}^T \boldsymbol{D}_\star \boldsymbol{b}_{\mathrm{reg}},
\label{eq:cc:sigma_ygivenxz}
\end{align}
with the regression coefficient vector
\begin{equation}
\boldsymbol{b}_{\mathrm{reg}} = \boldsymbol{D}_\star^{-1}
\begin{pmatrix}0\\[2pt]A_{12}\end{pmatrix},
\qquad
\boldsymbol{D}_\star =
\begin{pmatrix}
  \sigma_{1,\star}^2 & -\delta\sigma_{\mathrm{int},x}^2\\[4pt]
  -\delta\sigma_{\mathrm{int},x}^2 & A_{22} + \delta^2\sigma_{\mathrm{int},x}^2
\end{pmatrix}.
\label{eq:cc:b_D}
\end{equation}
The variance $\sigma_{y|\hat{x}\hat{z}}^2$ is independent of $y_{\mathrm{TF},\star}$.

\paragraph{Marginalizing $y_{\mathrm{TF},\star}$.}
The remaining integral over $y_{\mathrm{TF},\star}$ uses the joint posterior
$p(y_{\mathrm{TF},\star} \mid \hat{x}_\star, \hat{z}_\star, \theta)$, which
exploits the information that $\hat{z}_\star$ carries about $y_{\mathrm{TF},\star}$
through the dependence of $\hat{z}_\star$ on $y_{\mathrm{TF},\star}$.
Marginalizing the trivariate Eq.~\eqref{eq:cc:trivariate} over $\hat{y}_\star$
leaves a bivariate Gaussian likelihood
\begin{equation}
\begin{pmatrix}\hat{x}_\star\\ \hat{z}_\star\end{pmatrix}
\;\Big|\;y_{\mathrm{TF},\star},\theta
\;\sim\;
\mathcal{N}_2\!\bigl(\boldsymbol{m}_\star(y_{\mathrm{TF},\star}),\,\boldsymbol{D}_\star\bigr),
\label{eq:cc:bivariate_xz}
\end{equation}
with covariance $\boldsymbol{D}_\star$ as in Eq.~\eqref{eq:cc:b_D} and mean
\begin{equation}
\boldsymbol{m}_\star(y_{\mathrm{TF},\star})
= \boldsymbol{m}_\star(0) + y_{\mathrm{TF},\star}\,\boldsymbol{b}_{xz},
\qquad
\boldsymbol{b}_{xz} = \begin{pmatrix}1/\alpha\\[2pt] 1 - \delta/\alpha\end{pmatrix}.
\label{eq:cc:b_xz}
\end{equation}
Combined with the uniform prior on $y_{\mathrm{TF},\star}$, the posterior is a
truncated normal on $[y_{\min}, y_{\max}]$ whose untruncated precision and
mean are
\begin{equation}
\xi_{xz} = \boldsymbol{b}_{xz}^T \boldsymbol{D}_\star^{-1} \boldsymbol{b}_{xz},
\qquad
\mu^\dagger_{xz}
= \xi_{xz}^{-1}\,\boldsymbol{b}_{xz}^T \boldsymbol{D}_\star^{-1}
  \bigl(\boldsymbol{o}_\star - \boldsymbol{m}_\star(0)\bigr),
\label{eq:cc:T_post_xz}
\end{equation}
with $\boldsymbol{o}_\star = (\hat{x}_\star, \hat{z}_\star)^T$.
Let $\mu_{T|xz}(\theta)$ and $V_{T|xz}(\theta)$ denote the mean and variance of
this truncated normal (cf.\ Appendix~\ref{sec:infer_single_galaxy_y}).

Because $\mu_{y|\hat{x}\hat{z}}(y_{\mathrm{TF},\star})$ is linear in
$y_{\mathrm{TF},\star}$ with slope
\begin{equation}
s = 1 - \boldsymbol{b}_{\mathrm{reg}}^T \boldsymbol{b}_{xz},
\label{eq:cc:slope}
\end{equation}
the law of total expectation and variance give the exact conditional moments
\begin{align}
\mathrm{E}[\hat{y}_\star \mid \hat{x}_\star,\,\hat{z}_\star,\,\theta,\,\gamma]
&= \mu_{y|\hat{x}\hat{z}}\!\bigl(\mu_{T|xz}(\theta)\bigr),
\label{eq:cc:mean_ystar}\\[4pt]
\mathrm{Var}[\hat{y}_\star \mid \hat{x}_\star,\,\hat{z}_\star,\,\theta,\,\gamma]
&= \sigma_{y|\hat{x}\hat{z}}^2 + s^2\,V_{T|xz}(\theta).
\label{eq:cc:var_ystar}
\end{align}

\paragraph{Posterior predictive over MCMC draws.}
The posterior predictive mean and variance over draws
$\{(\theta^{(m)},\gamma^{(m)})\}_{m=1}^M$ are
\begin{align}
\overline{y}_\star
&= \frac{1}{M}\sum_{m=1}^M
   \mathrm{E}^{(m)}[\hat{y}_\star \mid \hat{x}_\star,\hat{z}_\star,\theta^{(m)},\gamma^{(m)}],
\label{eq:cc:pp_mean}\\[4pt]
\mathrm{Var}_{pp}(\hat{y}_\star)
&= \frac{1}{M}\sum_{m=1}^M
   \mathrm{Var}^{(m)} +
   \frac{1}{M}\sum_{m=1}^M
   \!\left(\mathrm{E}^{(m)} - \overline{y}_\star\right)^2.
\label{eq:cc:pp_var}
\end{align}

\paragraph{Prediction from rotation velocity alone.}
A reduced prediction that ignores $\hat{z}_\star$ is obtained by marginalizing
the trivariate Eq.~\eqref{eq:cc:trivariate} over $\hat{z}_\star$ rather than
conditioning on it.  The upper-left $2\times 2$ block of $\boldsymbol{B}_\star$
(Eq.~\eqref{eq:cc:B}) has $B_{12} = 0$, so $\hat{x}_\star$ and $\hat{y}_\star$
are independent at fixed $y_{\mathrm{TF},\star}$ and
\begin{equation}
\hat{y}_\star \mid y_{\mathrm{TF},\star}
\;\sim\;
\mathcal{N}\!\bigl(y_{\mathrm{TF},\star} + \Delta_{r,\star},\;A_{11}\bigr),
\label{eq:cc:y_given_yTF_xonly}
\end{equation}
with $A_{11} = \gamma^2\tau_c^2 + \sigma_{\mathrm{int},y}^2 + \sigma_{y,\star}^2$
as before.  The posterior of $y_{\mathrm{TF},\star}$ given $\hat{x}_\star$
alone is the truncated normal of the baseline model
(Eq.~\eqref{eq:post_yTF_tophat}), with mean $\mu_\star(\theta)$ and variance
$V_\star(\theta)$.  By the law of total expectation and variance,
\begin{align}
\mathrm{E}[\hat{y}_\star \mid \hat{x}_\star,\,\theta,\,\gamma]
&= \mu_\star(\theta) + \Delta_{r,\star},
\label{eq:cc:mean_xonly}\\[4pt]
\mathrm{Var}[\hat{y}_\star \mid \hat{x}_\star,\,\theta,\,\gamma]
&= A_{11} + V_\star(\theta).
\label{eq:cc:var_xonly}
\end{align}
Posterior-predictive moments over MCMC draws follow
Eqs.~\eqref{eq:cc:pp_mean}--\eqref{eq:cc:pp_var} with $\mathrm{E}^{(m)}$ and
$\mathrm{Var}^{(m)}$ replaced by the above.  This $x$-only prediction is
structurally identical to the baseline model
(Appendix~\ref{sec:infer_single_galaxy_y}) with $\sigma_{2,\star}^2$ replaced
by $A_{11}$.  Relative to the full $(\hat{x}_\star,\hat{z}_\star)$ prediction,
it discards the $y_{\mathrm{TF},\star}$ information carried by
$\hat{z}_\star$ but avoids propagating $z$-band photometric errors through
$\boldsymbol{b}_{\mathrm{reg}}$; which prediction is tighter depends on the
balance between these two effects.

\subsection{Reduction to Baseline}
\label{sec:cc:reduction}

Setting $\gamma = 0$ and $\alpha_{k,r} = \alpha_{k,z} = 0$ gives $A_{11} = \sigma_{\mathrm{int},y}^2 + \sigma_{y,i}^2 \equiv \sigma_{2,i}^2$,
$A_{12} = 0$, $\Delta_{r,i} = \Delta_{z,i} = 0$, and $\boldsymbol{b}_{\mathrm{reg}} = \boldsymbol{0}$.
The $\hat{z}$ measurement drops out of the prediction entirely:
$\sigma_{y|\hat{x}\hat{z}}^2 = A_{11} = \sigma_{2,i}^2$, recovering the
baseline model of Appendix~\ref{sec:likelihood}.
The likelihood at each galaxy reduces to the product
$\mathcal{N}(\hat x_i; (y_{\mathrm{TF},i}-\beta)/\alpha, \sigma_{1,i}^2)\,
 \mathcal{N}(\hat y_i; y_{\mathrm{TF},i}, \sigma_{2,i}^2)\,
 \mathcal{N}(\hat z_i; y_{\mathrm{TF},i} - \mu_c - \delta(\mu_{x,i}-\bar{x}), \sigma_{\mathrm{int},z}^2 + \sigma_{z,i}^2)$,
where the $\hat{z}$ factor is independent of $\hat{y}$ and contributes only
to the estimation of $(\mu_c, \delta, \sigma_{\mathrm{int},z})$, not to the
$r$-band prediction.

\subsection{Model Limitations}
\label{sec:cc:limitations}

The population color distribution is modeled as a single Gaussian
$\mathcal{N}(\mu_c + \delta(x_i-\bar{x}),\tau_c^2)$.
Non-Gaussian tails from dust, AGN contamination, or galaxy-type mixing are not
captured; a mixture or nonparametric prior on $c_i$ is a natural extension.
The independence of $\eta_i$ from $x_i$ (Eq.~\ref{eq:cc:eta_def}) is a model
assumption equivalent to Gaussianity of the color--velocity residuals; it should
be verified empirically against the observed $(\hat{c}_i, \hat{x}_i)$ scatter
before the model is applied.
The hyperparameter $\tau_c$ controls how much the $z$-band observation informs
the latent $y_i$; naive estimation of $\tau_c$ from the observed color variance
overestimates the intrinsic scatter because $\hat{c}_i = \hat{y}_i - \hat{z}_i$
has measurement noise $\sigma_{y,i}^2 + \sigma_{z,i}^2$; the posterior on $\tau_c$
from the full hierarchical fit automatically deconvolves this.

\section{Two-Color Extension}
\label{sec:2color}

The single-color model of \S\ref{sec:color_correction} uses the $r{-}z$ color
(one latent factor $\eta_i$) to sharpen the absolute-magnitude prediction.  A
natural extension adds $g$-band photometry, modeling the joint chromatic scatter
of the three absolute magnitudes directly.  This appendix describes the
\emph{2color} model, which replaces the trivariate likelihood
$(\hat{x}_i,\hat{y}_i,\hat{z}_i)$ with a quadrivariate likelihood
$(\hat{x}_i,\hat{y}_i,\hat{z}_i,\hat{g}_i)$.

\subsection{Generative Model}
\label{sec:2c:generative}

Each galaxy retains the latent TF variables $y_{\mathrm{TF},i}$ and $x_i$ from
\S\ref{sec:cc:generative}. Rather than modeling the $r{-}z$ and $g{-}r$ colors
through two \emph{independent} latent color factors, the 2color model represents
the joint intrinsic scatter of the three absolute magnitudes directly. Define
the mean absolute-magnitude relations for the $r$, $z$, and $g$ bands,
\begin{align}
  m_{y,i}  &= y_{\mathrm{TF},i} + \Delta_{r,i},\\
  m_{z,i}  &= y_{\mathrm{TF},i} - \mu_c - \delta_c(x_i-\bar{x}) + \Delta_{z,i},\\
  m_{g,i}  &= y_{\mathrm{TF},i} - \mu_g - \delta_g(x_i-\bar{x}) + \Delta_{g,i},
\end{align}
where $\mu_c,\mu_g$ are the mean $r{-}z$ and $g{-}r$ colors at $x=\bar{x}$,
$\delta_c,\delta_g$ the color--velocity slopes, and $\Delta_{r,i},\Delta_{z,i},
\Delta_{g,i}$ the band residual $k$-corrections
[$\Delta_{b,i}\equiv\alpha_{k,b}(\ln(1+z_{\mathrm{obs},i})-\overline{\ln(1+z)})$].
The latent absolute magnitudes $(y_i,z_i,g_i)$ scatter about these means with a
\emph{rank-2} $3\times3$ intrinsic covariance,
\begin{equation}
  \begin{pmatrix} y_i\\ z_i\\ g_i \end{pmatrix}
  \Bigm|\, y_{\mathrm{TF},i}, x_i
  \;\sim\;
  \mathcal{N}\!\left(
    \begin{pmatrix} m_{y,i}\\ m_{z,i}\\ m_{g,i} \end{pmatrix},\;
    \boldsymbol{S}
  \right),
  \qquad
  \boldsymbol{S} = \boldsymbol{V}\,\boldsymbol{\Sigma}_c\,\boldsymbol{V}^T,
  \label{eq:2c:S}
\end{equation}
where $\boldsymbol{V}$ is a $3\times2$ matrix with orthonormal columns spanning
the plane orthogonal to a \emph{free} unit null direction $\hat{\boldsymbol{n}}$,
and $\boldsymbol{\Sigma}_c$ is a free $2\times2$ symmetric positive-definite
matrix.
[There should be a mention that runs with a rank-3 covariance 
showed near degeneracy, which led to slow convergence.  This is 
why we go with rank-2.  Could this be true since \sigma_{x,int}
is another free parameter?]
  By construction $\boldsymbol{S}\,\hat{\boldsymbol{n}}=\boldsymbol{0}$:
$\boldsymbol{S}$ is rank~2, with a null space along $\hat{\boldsymbol{n}}$, so
exactly one direction in $(y,z,g)$ space carries no intrinsic scatter --- but
\emph{which} direction is inferred from the data, not fixed in advance.  The unit
vector $\hat{\boldsymbol{n}}$ is assigned a uniform prior on the sphere, and
$\boldsymbol{V}$ is constructed as an orthonormal basis of its orthogonal
complement (via a Householder reflector, so the map is smooth in
$\hat{\boldsymbol{n}}$).
The $2\times2$ covariance is parameterized as
$\boldsymbol{\Sigma}_c = \boldsymbol{D}_c\,\boldsymbol{R}_c\,\boldsymbol{D}_c$
with $\boldsymbol{D}_c=\mathrm{diag}(S_{c,1},S_{c,2})$ (two scatter
scales) and $\boldsymbol{R}_c$ a $2\times2$ correlation matrix (one free
parameter).  The intrinsic scatter therefore has 5 free parameters in total ---
two for $\hat{\boldsymbol{n}}$, two scales, and one correlation --- which is
exactly the dimension of a general rank-2 positive-semidefinite $3\times3$
matrix; it sits strictly between the fixed-null chromatic-only form (3
parameters) and an unconstrained full-rank covariance (6).  Only the plane
$\mathrm{span}(\boldsymbol{V})$, equivalently $\hat{\boldsymbol{n}}$ up to sign,
is identified: the in-plane orientation of the columns of $\boldsymbol{V}$ is a
gauge freedom absorbed by the full $2\times2$ $\boldsymbol{\Sigma}_c$.
Writing $S_{yy},S_{yz},S_{yg},S_{zz},S_{zg},S_{gg}$ for the entries of
$\boldsymbol{S}=\boldsymbol{V}\boldsymbol{\Sigma}_c\boldsymbol{V}^T$,
the measurement equations are
$\hat{y}_i\mid y_i\sim\mathcal{N}(y_i,\sigma_{y,i}^2)$ and likewise for
$\hat{z}_i,\hat{g}_i$, and the intrinsic (measurement-free) covariance of
$(y_i,z_i,g_i)$ enters the likelihood through the $\boldsymbol{B}_i$ matrix below.

\subsection{Marginalization to Quadrivariate Likelihood}
\label{sec:2c:marginalization}

Marginalizing $x_i$ over its intrinsic scatter $\sigma_{\mathrm{int},x}$ gives
a quadrivariate Gaussian for
$(\hat{x}_i,\hat{y}_i,\hat{z}_i,\hat{g}_i)\mid y_{\mathrm{TF},i}$
with mean vector
\begin{equation}
\boldsymbol{\mu}_i(y_{\mathrm{TF},i})
=
\begin{pmatrix}
  \mu_{x,i}(y_{\mathrm{TF},i})\\[2pt]
  y_{\mathrm{TF},i} + \Delta_{r,i}\\[2pt]
  y_{\mathrm{TF},i} + \Delta_{z,i} - \mu_c - \delta_c(\mu_{x,i}-\bar{x})\\[2pt]
  y_{\mathrm{TF},i} + \Delta_{g,i} - \mu_g - \delta_g(\mu_{x,i}-\bar{x})
\end{pmatrix}
\label{eq:2c:mean}
\end{equation}
and $4\times4$ covariance matrix
\begin{equation}
\boldsymbol{B}_i =
\begin{pmatrix}
  \sigma_{1,i}^2 & 0
    & -\delta_c\sigma_{\mathrm{int},x}^2
    & -\delta_g\sigma_{\mathrm{int},x}^2 \\[4pt]
  0 & A_{11}
    & A_{12}
    & A_{14} \\[4pt]
  -\delta_c\sigma_{\mathrm{int},x}^2 & A_{12}
    & A_{22} + \delta_c^2\sigma_{\mathrm{int},x}^2
    & A_{zg} \\[4pt]
  -\delta_g\sigma_{\mathrm{int},x}^2 & A_{14}
    & A_{zg}
    & A_{44} + \delta_g^2\sigma_{\mathrm{int},x}^2
\end{pmatrix},
\label{eq:2c:B}
\end{equation}
where the entries are set directly by the free intrinsic covariance
$\boldsymbol{S}$ (Eq.~\ref{eq:2c:S}),
\begin{align}
  A_{11} &= S_{yy} + \sigma_{y,i}^2, \label{eq:2c:A11}\\
  A_{12} &= S_{yz},\\
  A_{14} &= S_{yg},\\
  A_{22} &= S_{zz} + \sigma_{z,i}^2,\\
  A_{44} &= S_{gg} + \sigma_{g,i}^2,\\
  A_{zg} &= S_{zg} + \delta_c\delta_g\sigma_{\mathrm{int},x}^2. \label{eq:2c:A44}
\end{align}
[In the following there is no need to describe this in the context of the earlier parameterization.]
The $(3,4)$ entry has two contributions. $S_{zg}$ is a \emph{free} parameter [is this true in this parameterization?]: the
intrinsic $\hat{z}$--$\hat{g}$ covariance is no longer forced to zero, so
$\hat{z}$ and $\hat{g}$ are not required to be conditionally independent given
$(y_{\mathrm{TF},i}, x_i)$. Fixing $S_{zg}=0$ recovers the earlier
two-independent-latent-color-factor model. The second term is \emph{induced} by
the marginalization rather than free: because $m_{z,i}$ and $m_{g,i}$ both depend
on $x_i$ with slopes $-\delta_c$ and $-\delta_g$
(Eqs.~\ref{eq:2c:S}~ff.), integrating $x_i$ over its intrinsic scatter
correlates $\hat{z}$ and $\hat{g}$ by
$(-\delta_c)(-\delta_g)\sigma_{\mathrm{int},x}^2$, exactly as the same
marginalization produces the $-\delta_c\sigma_{\mathrm{int},x}^2$ entries in the
first row and column and the $\delta_c^2\sigma_{\mathrm{int},x}^2$,
$\delta_g^2\sigma_{\mathrm{int},x}^2$ terms on the diagonal. Retaining it is what makes $\boldsymbol{B}_i$ the marginal covariance of the
generative model, and hence positive definite by construction: $\boldsymbol{S}$
is positive semi-definite (rank~2, null along $\hat{\boldsymbol{n}}$) and the
strictly positive per-band measurement variances
$(\sigma_{y,i}^2,\sigma_{z,i}^2,\sigma_{g,i}^2)$ fill the null direction,
guaranteeing that $\boldsymbol{B}_i$ is strictly positive definite.  Dropping the
induced term admits parameter values for which $\boldsymbol{B}_i$ has a negative
eigenvalue and the Cholesky factorization fails.

\subsection{Closed-Form Likelihood}
\label{sec:2c:likelihood}

The integral over $y_{\mathrm{TF},i}$ proceeds exactly as in
\S\ref{sec:cc:marginalization}: the mean $\boldsymbol{\mu}_i$ is affine in
$y_{\mathrm{TF},i}$ with direction vector
\begin{equation}
\boldsymbol{b} =
\begin{pmatrix}
  1/\alpha\\[2pt] 1\\[2pt] 1-\delta_c/\alpha\\[2pt] 1-\delta_g/\alpha
\end{pmatrix},
\label{eq:2c:bvec}
\end{equation}
so the closed-form log-likelihood is
\begin{equation}
\ln\mathcal{L}_i
= -\tfrac{3}{2}\ln(2\pi)
  - \sum_k \ln L_{kk}
  - \tfrac{1}{2}\ln\xi_i
  - \tfrac{1}{2}Q_{0,i}
  + \ln\Delta\Phi_i
  - \ln(y_{\max}-y_{\min}),
\label{eq:2c:loglik}
\end{equation}
where $L_{kk}$ are the diagonal elements of the Cholesky factor of
$\boldsymbol{B}_i$, and
\begin{align}
  \xi_i &= \boldsymbol{b}^T \boldsymbol{B}_i^{-1} \boldsymbol{b},\\
  \mu^\dagger_i &= \phi_i/\xi_i, \quad
  \phi_i = \boldsymbol{b}^T \boldsymbol{B}_i^{-1} \boldsymbol{d}_i,\\
  Q_{0,i} &= \boldsymbol{d}_i^T\boldsymbol{B}_i^{-1}\boldsymbol{d}_i
             - \phi_i^2/\xi_i,\\
  \Delta\Phi_i &= \Phi\!\bigl(\sqrt{\xi_i}(y_{\max}-\mu^\dagger_i)\bigr)
                 -\Phi\!\bigl(\sqrt{\xi_i}(y_{\min}-\mu^\dagger_i)\bigr),
\end{align}
with $\boldsymbol{d}_i = \mathrm{obs}_i - \boldsymbol{\mu}_i(0)$.  The
normalization uses $-\frac{3}{2}\ln(2\pi)$ (rather than $-\frac{1}{2}\ln(2\pi)$
for the trivariate model) because the quadrivariate Gaussian has been reduced to
a 1D integral; i.e.,\ $4-1=3$ Gaussian dimensions are integrated analytically.

\subsection{Posterior Predictive Distribution}
\label{sec:2c:predict}

Prediction proceeds as in \S\ref{sec:cc:predict} but with a $3\times 3$
conditioning matrix $\boldsymbol{D}_\star$ formed from the $(\hat{x},\hat{z},\hat{g})$
sub-block of $\boldsymbol{B}_\star$:
\begin{equation}
\boldsymbol{D}_\star =
\begin{pmatrix}
  \sigma_{1,\star}^2
    & -\delta_c\sigma_{\mathrm{int},x}^2
    & -\delta_g\sigma_{\mathrm{int},x}^2\\[4pt]
  -\delta_c\sigma_{\mathrm{int},x}^2
    & A_{22}+\delta_c^2\sigma_{\mathrm{int},x}^2
    & A_{zg}\\[4pt]
  -\delta_g\sigma_{\mathrm{int},x}^2
    & A_{zg}
    & A_{44}+\delta_g^2\sigma_{\mathrm{int},x}^2
\end{pmatrix},
\label{eq:2c:D}
\end{equation}
with the $(\hat{z},\hat{g})$ entry $A_{zg}=S_{zg}+\delta_c\delta_g
\sigma_{\mathrm{int},x}^2$ inherited directly from $\boldsymbol{B}_\star$
(Eq.~\ref{eq:2c:A44}): both the free intrinsic $S_{zg}$ and the term induced by
marginalizing the shared latent $x_\star$ are retained, since
$\boldsymbol{D}_\star$ is a genuine sub-block of $\boldsymbol{B}_\star$. Dropping
the induced $\delta_c\delta_g\sigma_{\mathrm{int},x}^2$ term would make
$\boldsymbol{D}_\star$ inconsistent with $\boldsymbol{B}_\star$ and forfeit the
positive-definiteness it inherits from the latter.
The regression coefficient vector is
$\boldsymbol{b}_{\mathrm{reg}} = \boldsymbol{D}_\star^{-1}(0, A_{12}, A_{14})^T$
and the conditional moments are
\begin{align}
  \mu_{y|\hat{x}\hat{z}\hat{g}}(y_{\mathrm{TF},\star})
  &= y_{\mathrm{TF},\star} + \Delta_{r,\star}
    + \boldsymbol{b}_{\mathrm{reg}}^T\boldsymbol{r}_\star(y_{\mathrm{TF},\star}),
  \label{eq:2c:mu_cond}\\
  \sigma_{y|\hat{x}\hat{z}\hat{g}}^2
  &= A_{11}
    - (0,A_{12},A_{14})\,\boldsymbol{D}_\star^{-1}\,(0,A_{12},A_{14})^T,
  \label{eq:2c:sigma_cond}
\end{align}
where $\boldsymbol{r}_\star$ is the 3-vector of residuals
$(\hat{x}_\star - \mu_{x,\star},\;
  \hat{z}_\star - m_{z,\star},\;
  \hat{g}_\star - m_{g,\star})$
evaluated at $y_{\mathrm{TF},\star}$.

The posterior on $y_{\mathrm{TF},\star}$ given $(\hat{x}_\star,\hat{z}_\star,\hat{g}_\star)$ is a truncated normal with precision
\begin{equation}
  \xi_{xzg}
  = \boldsymbol{b}_{xzg}^T \boldsymbol{D}_\star^{-1} \boldsymbol{b}_{xzg},
  \qquad
  \boldsymbol{b}_{xzg}
  = \begin{pmatrix}1/\alpha\\[2pt]1-\delta_c/\alpha\\[2pt]1-\delta_g/\alpha\end{pmatrix},
  \label{eq:2c:xi_xzg}
\end{equation}
and mean $\mu^\dagger_{xzg} = \xi_{xzg}^{-1}\,\boldsymbol{b}_{xzg}^T
\boldsymbol{D}_\star^{-1}(\boldsymbol{o}_\star - \boldsymbol{m}_\star(0))$.
The exact posterior-predictive moments follow
Eqs.~\eqref{eq:cc:mean_ystar}--\eqref{eq:cc:pp_var} with
$\boldsymbol{D}_\star$, $\boldsymbol{b}_{\mathrm{reg}}$, $\xi_{xzg}$, and
$\boldsymbol{b}_{xzg}$ replacing their $2\times 2$ counterparts.

\subsection{Prediction from Rotation Velocity Alone}
\label{sec:2c:xonly}

Marginalizing $\hat{z}_\star$ and $\hat{g}_\star$ out of
Eq.~\eqref{eq:2c:B} gives
$\hat{y}_\star\mid y_{\mathrm{TF},\star}
\sim\mathcal{N}(y_{\mathrm{TF},\star}+\Delta_{r,\star},\,A_{11})$
(identical to Eq.~\eqref{eq:cc:y_given_yTF_xonly}), with
$A_{11}=S_{yy}+\sigma_{y,\star}^2$ (Eq.~\ref{eq:2c:A11}); only the $y$-band
intrinsic variance $S_{yy}$ enters, since $\hat{z}_\star$ and $\hat{g}_\star$
are marginalized out.
The prediction follows Eqs.~\eqref{eq:cc:mean_xonly}--\eqref{eq:cc:var_xonly}
with this~$A_{11}$.

\subsection{New Parameters and Priors}
\label{sec:2c:priors}

Relative to the single-color model, the 2color model adds the $g$-band mean
structure ($\delta_g,\mu_g,\alpha_{k,g}$) and replaces the $r{-}z$/$g{-}r$
product-form intrinsic-scatter parameters
($\gamma,\tau_c,\gamma_g,\tau_g,\sigma_{\mathrm{int},y},
\sigma_{\mathrm{int},z},\sigma_{\mathrm{int},g}$) with the rank-2
intrinsic covariance $\boldsymbol{S}$ of Eq.~\eqref{eq:2c:S}, parameterized by
the free null direction $\hat{\boldsymbol{n}}$, the two scatter scales
$S_{c,1},S_{c,2}$, and the $2\times2$ correlation matrix $\boldsymbol{R}_c$:

\begin{center}
\begin{tabular}{lll}
\hline
Parameter & Meaning & Prior \\
\hline
$\hat{\boldsymbol{n}}$ & rank-2 null direction & uniform on $S^2$ \\
$S_{c,1},S_{c,2}$ & scatter scales & $\mathrm{HalfCauchy}(0,1)$ \\
$\boldsymbol{R}_c$ & $2\times2$ correlation & $\mathrm{LKJ}(2)$ \\
$\delta_g$ & population $g{-}r$ color--velocity slope & $\mathcal{N}(0,1)$ \\
$\mu_g$ & mean $g{-}r$ color at $x=\bar{x}$ & $\mathcal{N}(\bar{c}_{g,\mathrm{obs}},1)$ \\
$\alpha_{k,g}$ & $g$-band $k$-correction slope & $\mathcal{N}(0,5)$ \\
\hline
\end{tabular}
\end{center}

The null direction $\hat{\boldsymbol{n}}$ is a unit vector with a uniform prior
on the sphere, and $\boldsymbol{R}_c$ is sampled through its Cholesky factor (a
single free element) with an $\mathrm{LKJ}(2)$ prior.  The total parameter count
for $\boldsymbol{S}$ is 5 (two for $\hat{\boldsymbol{n}}$, two scales, one
correlation), between the 3 of the fixed-null chromatic-only special case and the
6 of an unconstrained free $3\times3$ parameterization.  The rank-2 restriction
eliminates a near-degenerate ridge in the posterior that arises because
photometry alone cannot distinguish intrinsic achromatic scatter from
peculiar-velocity-induced scatter: with a free $3\times3$ $\boldsymbol{S}$, the
fitted matrix is dominated by a single near-achromatic mode ($\sim\!99\%$
correlated, $\sim\!0.4$~mag) that is degenerate with $C_{vv}$, degrading HMC
efficiency (mean treedepth at maximum, ESS~$\sim$\!0) and causing Cholesky
failures on $\boldsymbol{B}_i$ when the sampler explores the boundary of the
degenerate ridge.  Restricting $\boldsymbol{S}$ to rank~2 places the sampler on a
well-defined lower-dimensional manifold and removes this pathology, \emph{without}
hard-coding which direction is excluded: the null direction is fit, so the
achromatic axis $\boldsymbol{e}$ is no longer imposed but can be compared against
the posterior of $\hat{\boldsymbol{n}}$.

\paragraph{Achromatic variance and the downstream velocity model.}
The three absolute magnitudes share a common distance modulus: any
peculiar-velocity error shifts $(y_i,z_i,g_i)$ by an identical amount along the
achromatic direction $\boldsymbol{e}=(1,1,1)/\!\sqrt{3}$.  The rank-2 model
excludes exactly one direction from the intrinsic scatter
($\boldsymbol{S}\,\hat{\boldsymbol{n}}=\boldsymbol{0}$; Eq.~\ref{eq:2c:S}), that
direction being inferred.  If the fitted $\hat{\boldsymbol{n}}$ coincides with
$\boldsymbol{e}$, the excluded direction is precisely the
peculiar-velocity/distance-modulus mode: the calibration model then contains
\emph{no} peculiar-velocity scatter, all achromatic variance is owned by the
downstream velocity-covariance model $C_{vv}$ (which adds a coherent-flow prior to
the per-galaxy distance-modulus posteriors of the prediction step,
\S\ref{sec:2c:predict}), and the separation is clean --- $\boldsymbol{S}$ in the
chromatic plane $\perp\boldsymbol{e}$, $C_{vv}$ along $\boldsymbol{e}$, no
double-counting.  We report the posterior angle between $\hat{\boldsymbol{n}}$ and
$\boldsymbol{e}$ as a diagnostic; a posterior concentrated near $0^\circ$
\emph{measures} the alignment that the fixed-null chromatic-only model would
otherwise assume.

This design targets a specific risk: achromatic intrinsic scatter is only weakly
identifiable from photometry alone --- it is degenerate with
peculiar-velocity-induced scatter --- so fitting a full-rank $\boldsymbol{S}$
lets an achromatic component absorb an unknown fraction of the PV signal and bias
the downstream flow-field estimate.  The chromatic structure (intrinsic color
scatter, color--velocity slope residuals), by contrast, is fully identifiable
from the multi-band photometry and is precisely what the $g$-band extension is
designed to measure.  Because the excluded direction is weakly constrained, the
posterior on $\hat{\boldsymbol{n}}$ may be broad; that is itself the honest
outcome, and a posterior-predictive check on the color residuals ($z{-}y$,
$g{-}y$) validates the chromatic structure independently of it.

\subsection{Relation to the Single-Color Model}
\label{sec:2c:reduction}

The rank-2 model nests the fixed-null chromatic-only model as the special case
$\hat{\boldsymbol{n}}=\boldsymbol{e}=(1,1,1)/\!\sqrt{3}$, but for a generic fitted
null direction it does \emph{not} reduce to the single-color model of
\S\ref{sec:cc:marginalization}.  The single-color model carries independent
per-band intrinsic variances
($\sigma_{\mathrm{int},y}^2,\sigma_{\mathrm{int},z}^2$) that include achromatic
scatter, whereas $\boldsymbol{S}=\boldsymbol{V}\boldsymbol{\Sigma}_c
\boldsymbol{V}^T$ is rank~2 and satisfies
$\boldsymbol{S}\,\hat{\boldsymbol{n}}=\boldsymbol{0}$; it therefore cannot
reproduce a full-rank (in particular, a generic diagonal) intrinsic covariance ---
a rank-2 matrix cannot be full rank, whatever $\hat{\boldsymbol{n}}$.  The rank-2
restriction is a deliberate modeling choice (\S\ref{sec:2c:priors}), not a
special-case limit of the single-color model.

Setting the second scale $S_{c,2}=0$ (so $\boldsymbol{\Sigma}_c$ is rank~1) makes
$\boldsymbol{S}$ rank~1, $\boldsymbol{S}=S_{c,1}^2\,\boldsymbol{v}\boldsymbol{v}^T$
for the surviving column $\boldsymbol{v}$ of $\boldsymbol{V}$ (a unit vector in the
plane $\perp\hat{\boldsymbol{n}}$); additionally taking $\delta_g=0$ and
$\alpha_{k,g}=0$ makes $\hat{g}$ uninformative.  This is the analog of the
single-color model with a single intrinsic-scatter degree of freedom, but it still
carries no scatter along $\hat{\boldsymbol{n}}$ --- the independent
$\sigma_{\mathrm{int},y}^2$ and $\sigma_{\mathrm{int},z}^2$ diagonal terms of the
single-color $\boldsymbol{B}_i$ have no counterpart here.


\section{Redshift-Dependent Rotational-Velocity Correction}
\label{sec:vcorr}

\subsection{Motivation}
\label{sec:vc:motivation}

Rotational velocities are measured at a fiber placed at $r_{\mathrm{fib}} = f\,R_{26}$,
where $R_{26}$ is the $r$-band semi-major axis of the
$SB = 26~\mathrm{mag}\,\mathrm{arcsec}^{-2}$ isophote and $f$ is the fiber
fraction; for the DR2 measurements $f = 0.4$, corresponding to the tabulated
$V_{0.4R_{26}}$. Because cosmological surface-brightness dimming scales as
$(1+z)^4$, an isophote defined at a fixed \emph{observed} surface brightness
corresponds to a redshift-dependent \emph{physical} radius. We therefore want
every velocity to refer to the same physical location, $f\,R_{26}$ evaluated at
a fiducial redshift $z_{\mathrm{fid}} = 0.05$. Left uncorrected, the drift of
$R_{26}$ with redshift moves the fiber to a different physical radius as a
function of $z$, imprinting a monotonic-in-$z$ systematic on the log-distance
ratios that biases the $f\sigma_8$ inference.

The correction has three steps: (i) the surface-brightness shift of the isophote
relative to $z_{\mathrm{fid}}$; (ii) the conversion of that shift to a fractional
radius shift through the galaxy light profile; and (iii) the conversion of the
radius shift to a velocity correction through the local slope of the rotation
curve.

\subsection{Surface-brightness shift of the isophote}
\label{sec:vc:dmu}

Relative to the fiducial redshift, the observed surface brightness of a fixed
physical isophote is shifted by
\begin{equation}
\delta\mu(z) = 10\,\log_{10}\!\frac{1+z}{1+z_{\mathrm{fid}}}
  + \bigl[\,k_\mu(z) - k_\mu(z_{\mathrm{fid}})\,\bigr],
\label{eq:vc:dmu}
\end{equation}
where the leading coefficient $10 = 2.5\times 4$ is cosmological dimming and
$k_\mu$ is the surface-brightness $k$-correction. In this work the SGA photometry
is $K$-corrected to $z_{\mathrm{fid}}$ upstream and no separable per-galaxy
$k_\mu$ is available, so we adopt $k_\mu(z) - k_\mu(z_{\mathrm{fid}}) \approx 0$
and retain only the dimming term (\S\ref{sec:vc:caveats}).

\subsection{From surface brightness to radius via the light profile}
\label{sec:vc:radius}

We convert $\delta\mu$ to a radial shift using the SGA-2020 curve-of-growth
fit \citep{SGA},
\begin{equation}
m(r) = m_{\mathrm{tot}} + m_0\,\ln\!\bigl[\,1 + \alpha_1 (r/r_0)^{-\alpha_2}\,\bigr],
\label{eq:vc:cog}
\end{equation}
with the scale radius fixed at $r_0 = 10''$. The surface brightness is the
derivative of the enclosed flux per unit area; for elliptical isophotes of fixed
axis ratio $q$ the enclosed area is $A(r) = \pi q\, r^2$, so with
$F(<r) = 10^{-0.4\,m(r)}$ the local surface brightness is
\begin{equation}
\mu(r) = m(r) + 2.5\,\log_{10} r - 2.5\,\log_{10}\!\bigl[-m'(r)\bigr] + \mathrm{const},
\label{eq:vc:mu}
\end{equation}
where the per-galaxy constant (axis ratio and zeropoints) is independent of $r$.
The corrected isophote radius $r' = (1+\delta)\,R_{26}$ is the solution of
\begin{equation}
\mu(r') - \mu(R_{26}) = \delta\mu ,
\label{eq:vc:solve}
\end{equation}
which we evaluate directly from Eqs.~\eqref{eq:vc:cog}--\eqref{eq:vc:mu}. Because
the difference depends only on the ratio
$u(r)/u(R_{26}) = (r/R_{26})^{-\alpha_2}$, with
$u \equiv \alpha_1 (r/r_0)^{-\alpha_2}$, the scale radius $r_0$ and the
normalization $\alpha_1$ cancel, and $u(R_{26})$ is fixed by the measured
aperture magnitude. The fractional shift is
$\delta \equiv \delta r / R_{26} = r'/R_{26} - 1$; linearizing
Eq.~\eqref{eq:vc:solve} gives the closed form
\begin{equation}
\delta \simeq \frac{\delta\mu}{R_{26}\,\bigl(d\mu/dr\bigr)\big|_{R_{26}}} .
\label{eq:vc:dlin}
\end{equation}
For $z > z_{\mathrm{fid}}$ we have $\delta\mu > 0$: the dimmed isophote is observed
at smaller radius, so the correction moves the target radius \emph{outward},
$\delta > 0$.

\subsection{From radius to velocity}
\label{sec:vc:velocity}

A fractional change $\delta$ in the sampling radius maps to a velocity change
through the local logarithmic slope of the rotation curve
\citep{Yegorova2007},
\begin{equation}
V_{\mathrm{corr}} = V_{\mathrm{meas}}\,
  \Bigl[\,1 + \delta\,\frac{d\ln V}{d\ln r}\,\Bigr].
\label{eq:vc:vcorr}
\end{equation}
Adopting the \citet{Ravi2023} rotation-curve parameterization
$V(r) = V_{\max}\,r\,/\,(R_{\mathrm{turn}}^{\alpha_{\mathrm{RC}}}
+ r^{\alpha_{\mathrm{RC}}})^{1/\alpha_{\mathrm{RC}}}$,
\begin{equation}
\frac{d\ln V}{d\ln r} = 1
  - \frac{r^{\alpha_{\mathrm{RC}}}}{R_{\mathrm{turn}}^{\alpha_{\mathrm{RC}}}
  + r^{\alpha_{\mathrm{RC}}}} ,
\label{eq:vc:logslope}
\end{equation}
where we hold the shape parameter fixed at $\alpha_{\mathrm{RC}} = 2.5$ as a
provisional choice, carried as a symbol so the value is easy to revise. At the
fiber $r = f\,R_{26}$ the ratio $r/R_{\mathrm{turn}} = f/(R_{\mathrm{turn}}/R_{26})$
is independent of $R_{26}$, so Eq.~\eqref{eq:vc:logslope} reduces to a single
population value set by $R_{\mathrm{turn}}/R_{26}$.

\subsection{Linearized correction and typical magnitude}
\label{sec:vc:magnitude}

Combining Eqs.~\eqref{eq:vc:dmu}, \eqref{eq:vc:dlin}, and \eqref{eq:vc:vcorr},
\begin{equation}
\frac{V_{\mathrm{corr}}}{V_{\mathrm{meas}}} - 1
  \simeq \frac{d\ln V}{d\ln r}\;
  \frac{10\,\log_{10}\!\bigl[(1+z)/(1+z_{\mathrm{fid}})\bigr]}
       {R_{26}\,\bigl(d\mu/dr\bigr)\big|_{R_{26}}} .
\label{eq:vc:final}
\end{equation}
For the DR2 sample the SGA light profile gives a median inverse slope
$\bigl(R_{26}\,d\mu/dr\bigr)^{-1} \approx 0.19~\mathrm{mag}^{-1}$, so
$\delta \approx 0.076$ at $z = 0.15$. With the provisional
$R_{\mathrm{turn}}/R_{26} = 0.25$ (giving $d\ln V/d\ln r \approx 0.24$), the
velocity correction is $\approx +1.8\%$ at $z = 0.15$ relative to
$z_{\mathrm{fid}} = 0.05$, rising to $\approx +3\%$ near $z \approx 0.25$ and
vanishing by construction at $z = z_{\mathrm{fid}}$. The correction is small per
galaxy but systematic in redshift.

\subsection{Caveats}
\label{sec:vc:caveats}

Several approximations enter the correction. (i)~The log-slope $d\ln V/d\ln r$ is
a population average; the \citet{Ravi2023} MaNGA-matched rotation-curve fits are
not yet incorporated, so the ratio $R_{\mathrm{turn}}/R_{26}$ (and hence the
amplitude) is provisional and the correction scales linearly with it.
(ii)~The shape parameter $\alpha_{\mathrm{RC}} = 2.5$ is provisional.
(iii)~The radial conversion depends on the assumed SGA curve-of-growth light
profile. The SGA $R_{26}$ isophotes are defined at a fixed
\emph{observed}-frame surface brightness, so the dimming correction is warranted;
were they instead rest-frame or dimming-corrected the correction would not apply.
(iv)~The surface-brightness $k$-correction term $k_\mu$ is neglected. Although each per-galaxy correction is at the percent
level, its systematic dependence on redshift makes it relevant for $f\sigma_8$
science.


\bibliographystyle{aasjournal}

\bibliography{references}% Produces the bibliography via BibTeX.
\end{document}

%
% ****** End of file main.tex ******
